\chapter{The 24 Types of Semilinear Representation}\label{chapter:twelvefoldway}

In this chapter we conduct an in-depth study of the irreducible objects of $\Rep_{\bbL}^{\hspace{0.1em} \sigma}(G)$ in the case that the extension $\bL / \bK$ has degree $2$. We will associate to each $W \in \Irr_{\bbL}^{\hspace{0.1em} \sigma}(G)$ a \emph{type}, which characterises the behaviour of $W$ under the induction and restriction functors of diagram (\ref{eqn:diamond}), and in many cases allows the semilinear Schur index $m(W)$ of $W$ (cf.\ Theorem \ref{thm:facts}) to be determined by purely linear representation theoretic data. We make the following standing assumption.

\vspace{0.5em}
\begin{center}
\emph{We assume in this chapter that $\bL / \bK$ has degree $2$.}
\end{center}
\vspace{0.5em}

In particular, this assumption implies that $H \leq G$ is a Real structure on $H$ (cf.\ Chapter \ref{chapter:realstructures}). We will additionally assume that $\bK$ has characteristic two from Section \ref{sect:diamondanalysis} onwards.
 
 \subsection{Notation}\label{sect:chap3notation} Because $\bL / \bK$ is quadratic, $\sigma \colon G \twoheadrightarrow \Gamma = \Gal(\bL / \bK)$ is determined by the requirement that $H$ act trivially. We thus suppress $\sigma$ from the notation and write:
 \begin{itemize}
    \item $\Rep_{\bbL}^{\hspace{0.1em} \star}(G) \coloneqq \Rep_{\bbL}^{\hspace{0.1em} \sigma}(G)$.
    \item $\Irr_{\bbL}^{\hspace{0.1em} \star}(G) \coloneqq \Irr_{\bbL}^{\hspace{0.1em} \sigma}(G)$.
    \item $\bL_{\star \hspace{-0.1em}}[G] \coloneqq \bL_{\sigma} \hspace{-0.1em}[G]$.
    \item $\overline{(-)} \colon \bL \rightarrow \bL$ is the unique non-trivial element of $\Gamma$.
    \item For $V \in \Vect_{\bbL}$, $\overline{V} \in \Vect_{\bbL}$ is $V$ but with $\bL$-vector space structure by
    \[
    \alpha * v \coloneqq \overline{\alpha} v \ \mbox{  for all  } \ \alpha\in\bL , v\in V.
    \]
 \end{itemize}

 \subsection{Action of $C_2$ on $\Irr_{\bbL}(H)$} \label{sect:C2actionIrrH}
 Recall that we write $C_2 = \langle \epsilon \rangle$ for the group with two elements. There is an action of $C_2$ on $\Irr_{\bbL} (H)$,
 defined on $V \in \Irr_{\bbL}(H)$ by setting $\epsilon * V \coloneqq {}^{\omega}\overline{V}$, which is $\overline{V}$
 as a $\bL$-vector space but with action of $H$ defined by picking an element $\omega \in G \setminus H$, and setting 
 \[
 h * v \coloneqq (\omega^{-1} h \omega)(v)
 \]
 for any $h \in H$ and $v \in V$. This gives ${}^{\omega}\overline{V}$ the structure of a $\bL[H]$-module,
 which is independent (up to isomorphism) of the choice of $\omega$. Hence it defines an action of $C_2$ on the set $\Irr_{\bbL}(H)$.

 \subsection{The three S-types of irreducible representation}

 Under our assumption that $\bL / \bK$ has degree $2$, Theorem \ref{thm:facts} specialises to the following.

 \begin{thm}\label{thm:threeoptions}
    Suppose that $W \in \Irr_{\bbL}^{\hspace{0.1em} \star}(G)$ and $V \in \Irr_{\bbL}(H)$ are such that $V$ is an irreducible submodule of $W|_H$, or equivalently that $W$ is an irreducible submodule of $\Ind_H^G(V)$. Then $V$ and $W$ are in either case \emph{A}, \emph{B} or \emph{C} of Table \ref{tab:thewaysRHSsquare}.
 \end{thm}
 \begin{table}[h]
\caption{Induction and Restriction Possibilities Between $\Irr_{\bbL}(H)$ and $\Irr_{\bbL}^{\hspace{0.1em} \star}(G)$}
\label{tab:thewaysRHSsquare}
\begin{center}
\def\arraystretch{1.1}
\begin{tabular}{ c c c c }
    \toprule 
Case & \textbf{A} & \textbf{B} & \textbf{C} \\
\midrule
 $W|_H$ & $V$ & $V \oplus {}^{\omega}\overline{V}$ & $V \oplus V$ \\ 
 $\Ind_H^G(V)$ & $W \oplus W$ & $W$ & $W$\\ 
 $V \cong  {}^{\omega}\overline{V}$? & Yes & No & Yes \\
  $m(W)$ & 1 & 1 & 2 \\
  $\dim_{\bL} \End_{\bL[H]}(V) : \dim_{\bK} \End_{\bL_{\star \hspace{-0.1em}}[G]}(W)$ & $1:1$ & $1:2$ & $1:4$ \\ 
  \bottomrule
\end{tabular}
\end{center} 
\end{table}
 \begin{proof}
    Because $\bL / \bK$ is finite we may apply Theorem \ref{thm:facts}, noting that the action of $C_2$-defined above agrees with that of $\Gamma$ in Theorem \ref{thm:facts}. The $C_2$-orbit of $V$ has size either $1$ or $2$. If $V \not\cong  {}^{\omega}\overline{V}$, then $\Gamma_V = 1$, and so by Theorem \ref{thm:facts} $m(W) = 1$ and
    \[
        W|_H \cong V \oplus  {}^{\omega}\overline{V}, \qquad \Ind_{H}^G(V) \cong W.
    \]
    Furthermore,
    \begin{align*}
    \dim_{\bK} \End_{\bL_{\star \hspace{-0.1em}}[G]}(W) &= \dim_{\bL} \End_{\bL[H]}(V \oplus {}^{\omega}\overline{V}), \\
    &= \dim_{\bL} \End_{\bL[H]}(V) +  \dim_{\bL} \End_{\bL[H]}({}^{\omega}\overline{V}), \\
    &= 2 \dim_{\bL} \End_{\bL[H]}(V).
    \end{align*}
    If instead $\Gamma_V = \Gamma$, then as $m(W) | |\Gamma_V| = 2$, there are two options: $m(W) = 1$ or $m(W) = 2$. Similarly to the above, these correspond to the two remaining cases (A) and (C).
 \end{proof}

In particular, we see that Case (B) is distinguished from Cases (A) and (C) from the action of $C_2$ on $\Irr_{\bbL}(H)$. However this action doesn't differentiate between cases (A) and (C).

The main aim of this chapter is to improve on this and obtain a finer classification of $\Irr_{\bbL}^{\hspace{0.1em} \star}(G)$, by considering the other linear representation categories of Diagram \eqref{eqn:diamond} beyond $\Rep_{\bbL}(H)$.

\section{Index Two Induction and Restriction}\label{sect:fourfoldway}

We first consider general induction and restriction behaviours between algebras equipped with certain endofunctors on their categories of representations, which we will apply to the extensions of Diagram \eqref{eqn:diamondrings}. This will give us a way to establish versions of Theorem \ref{thm:threeoptions} for all other extensions of diagram (\ref{eqn:diamondrings}) under our standing assumption that $\bL / \bK$ has degree $2$.

Suppose in this section that $\bK$ is any field, $A$ and $B$ are $\bK$-algebras, and $f \colon A \rightarrow B$ is a $\bK$-algebra homomorphism. The restriction functor $\Res \colon \Mod_B \rightarrow \Mod_A$
has left and right adjoint functors: induction and coinduction 
\[
\Ind \coloneqq B\otimes_{A}(-), \ \Cnd \coloneqq \Hom_A(B,-) \colon \Mod_A \rightarrow \Mod_B \, .
\]
We assume existence of covariant $\bK$-linear dualities
\begin{align*}
    (-)^* &\colon \Mod_A \rightarrow \Mod_A, \\
    \overline{(-)} \hspace{0.4em} &\colon \Mod_B \rightarrow \Mod_B ,
\end{align*}
meaning covariant functors whose composition with themselves are naturally isomorphic to the identity functor, which are related in the sense that there are isomorphisms
\begin{align*}
    \Res(\Cnd(M)) &\cong M \oplus M^*, \\
    \Cnd(\Res(N)) &\cong N \oplus \overline{N}, \\
    \Cnd(M) &\cong \Cnd(M^*), \\
    \Res(N) &\cong \Res(\overline{N}).
\end{align*}
for all $M \in \Mod_A$, $N \in \Mod_B$.
If  $M_1, M_2 \in \Rep (A)$ and $N_1, N_2 \in \Rep (B)$, where $\Rep (A)$ denotes the full subcategory of modules finite-dimensional over $\bK$, and similarly for $B$, we set
\begin{align*}
    (M_1, M_2)_A &\coloneqq \dim_{\bbK} \Hom_A(M_1, M_2), \\
    (N_1, N_2)_B &\coloneqq \dim_{\bbK} \Hom_B(N_1, N_2).
\end{align*}
These distribute over direct sums in each variable and, because $\Res$ is left adjoint to $\Cnd$, satisfy the identity.
\begin{equation} \label{ResCnd}
    (\Res(N), M)_A = (N, \Cnd(M))_B.
\end{equation}
In particular, using~\eqref{ResCnd} and the properties above, we arrive at the  identities
\begin{align}
    (\Res(N), \Res(N))_A &= (N,N)_B + (N, \overline{N})_B, \label{eqn:innerprod1} \\
    (\Cnd(M), \Cnd(M))_B &= (M,M)_A + (M, M^*)_A. \label{eqn:innerprod2}
\end{align}

\begin{prop}\label{prop:algindressummary}
Suppose that $A$ and $B$ are as described at the start of this section, $M \in \Rep(A)$ and $N \in \Rep(B)$, finite-dimensional over $\bK$, and that $0 \neq (\Res(N),M)_A$, or equivalently, $0 \neq (N, \Cnd(M))_B$. Then there are four mutually exclusive possibilities, described in Table \ref{tab:indres}.

In particular, $\Res(N)$ and $\Cnd(M)$ are semisimple, the factors of $\Cnd(M)$ are a subset of $\{N, \overline{N}\}$, and the factors of $\Res(N)$ are a subset of $\{M, M^*\}$.
\end{prop}

\begin{table}
\caption{Induction and Restriction Possibilities Between $A$ and $B$}
\label{tab:indres}
\begin{center}
\def\arraystretch{1.1}
\begin{tabular}{ c c c c c } 
    \toprule
Case & \textbf{1} & \textbf{2} & \textbf{3(a)} & \textbf{3(b)} \\
\midrule
 $\Cnd(M)$ & $N \oplus \overline{N}$ & $N$ & $N \oplus N$ & $N$\\ 
 $\Res(N)$ & $M$ & $M \oplus M^*$ & $M$ &  $M \oplus M$ \\ 
 $M \cong M^*$? & Yes & No & Yes & Yes \\ 
 $N \cong \overline{N}$? & No & Yes & Yes & Yes \\ 
  $(M,M)_A : (N,N)_B$ & $1:1$ & $1:1$ & $2:1$ & $1:2$ \\ 
  \bottomrule
\end{tabular}
\end{center} 
\end{table}

\begin{proof}
Suppose that $M \in \Rep(A)$ and $N \in \Rep(B)$, so $M$ and $N$ are finite dimensional over $\bK$, and that $0 \neq (\Res(N),M)_A$, or equivalently, $0 \neq (N, \Cnd(M))_B$. We have the following cases.

\vspace{0.5em}
\textbf{Case 1:} $N \not\cong \overline{N}$. Then because $\Res(N) = \Res(\overline{N})$,
\begin{equation} \label{NoNCndM}
N \oplus \overline{N} \hookrightarrow \Cnd(M),
\end{equation}
and so because $\Res$ is faithful,
\[
\Res(N) \oplus \Res(N) \cong \Res(N \oplus \overline{N}) \hookrightarrow \Res(\Cnd(M)) \cong M \oplus M^*.
\]
Therefore $\Res(N) \cong M$ and $M \cong M^*$. Furthermore, the map in~\eqref{NoNCndM} must have been an isomorphism, and so $\Cnd(M) \cong N \oplus \overline{N}$. Comparing dimensions of endomorphism rings,
\begin{align*}
    2(M,M)_A &= (M,M)_A + (M,M^*)_A, \\
    &= (\Cnd(M), \Cnd(M))_B, \\
    &= (N \oplus \overline{N}, \Cnd(M))_B, \\
    &= (N,\Cnd(M))_B + (\overline{N}, \Cnd(M))_B, \\
    &= 2(N,\Cnd(M))_B, \\
    &= 2 (N,N)_B,
\end{align*}
hence $(M,M)_A = (N,N)_B$. Note that we have used the equalities
\[
(\overline{N}, \Cnd(M))_B = (\Res(\overline{N}), M)_A = (\Res(N), M)_A = (N, \Cnd(M)).
\]

\vspace{0.5em}
\textbf{Case 2:} $M \not\cong M^*$. We have that $N \hookrightarrow \Cnd(M)$, and therefore because $\Res$ is faithful,
\[
\Res(N) \hookrightarrow \Res(\Cnd(M)) \cong M \oplus M^*,
\]
and $\Res(N)$ is semisimple. Because $\Cnd(M) = \Cnd(M^*)$, $(\Res(N), M) = (\Res(N),M^*)$, hence
\[
M \oplus M^* \hookrightarrow \Res(N)
\]
because $\Res(N)$ is semisimple, and thus $\Res(N) \cong M \oplus M^*$. In particular,
\[
N \oplus \overline{N} \cong \Cnd(\Res(N)) \cong \Cnd(M) \oplus \Cnd(M^*),
\]
and $\Cnd(M) \cong N$, $N \cong \overline{N}$. Similarly to the previous case, $(M,M)_A = (N,N)_B$.

\vspace{0.5em}
\textbf{Case 3:} $N \cong \overline{N}$ and $M \cong M^*$. We have that $N \hookrightarrow \Cnd(M)$, and so, $\Res$ being faithful,
\[
\Res(N) \hookrightarrow \Res(\Cnd(M)) \cong M \oplus M.
\]
Therefore, there are two options.

\vspace{0.5em}
\textbf{Case 3(a):} $\Res(N) = M$. Then
\[
\Cnd(M) = \Cnd(\Res(N)) = N \oplus N.
\]
In this case, $(M,M)_A = 2 (N,N)_B$, as
\begin{align*}
    2(N,N)_B &= (N,N)_B + (N,\overline{N})_B, \\
    &= (\Res(N), \Res(N))_A, \\
    &= (M,M)_A.
\end{align*}

\vspace{0.5em}
\textbf{Case 3(b):} $\Res(N) = M \oplus M$. Then
\[
\Cnd(M) \oplus \Cnd(M) = \Cnd(\Res(N)) = N \oplus N,
\]
so $\Cnd(M) = N$. Similarly, $(N,N)_B = 2 (M,M)_A$.
\end{proof}

\section{Examples From the Diamond}\label{sect:examplesfromdiamond} Suppose from now we are in the context described at the start of this chapter, so $\bL / \bK$ is a Galois extension of degree two, and $H \leq G$ is a Real structure on $H$. We also fix an element $\omega \in G \setminus H$.

In this section we concentrate on understanding better the structure of the \emph{diamond}, by which we mean diagram (\ref{eqn:diamond}). More precisely, in this section we show that when $\charfield(\bK)\neq 2$, all six extensions of the diamond fit into the formalism developed in the previous section.

All six extensions of the diamond are Frobenius by Proposition \ref{prop:frobenius}, and so in what follows we consider $\Ind = B \otimes_A (-)$, rather than coinduction $\Cnd$ which we used above.

\subsection{The Extension $\bL[H] \hookrightarrow \bL_{\star \hspace{-0.1em}}[G]$}

Here the duality on $\Mod_{\bL_{\star \hspace{-0.1em}}[G]}$ is the identity functor, and the duality on $\Mod_{\bL[H]}$ sends $V \mapsto {}^{\omega}\overline{V}$. That this satisfies the required identities follows from \cite[Cor.\ 8.5]{TAYSchur}, \cite[Prop.\ 8.17]{TAYSchur} and Theorem \ref{thm:facts}. Because the duality on $\Mod_{\bL_{\star \hspace{-0.1em}}[G]}$ is trivial Case 1 doesn't occur. The remaining three cases recover the three cases of Table \ref{tab:thewaysRHSsquare}: Case 2 is Case (B), Case 3(a) is Case (A) and Case 3(b) is Case (C).

\subsection{The Extensions $\bK[H] \hookrightarrow \bL[H]$ and $\bK[G] \hookrightarrow \bL[G]$}\label{sect:classicalindextwofieldextoptions} Let us describe everything for $\bK[H] \hookrightarrow \bL[H]$, the case for $\bK[G] \hookrightarrow \bL[G]$ obtained by replacing $H$ by $G$ everywhere. This can be viewed as a special case of the previous example in light of \cite[\S10]{TAYSchur}.

For $\bK[H] \hookrightarrow \bL[H]$, the duality on $\Mod_{\bK[H]}$ is the identity functor, and on $\Mod_{\bL[H]}$ sends $V \mapsto \overline{V}$. Now it is the duality on $\Mod_{\bK[H]}$ that is trivial, so Case 2 doesn't occur.

We also have three cases arising from Theorem \ref{thm:facts}: if $X \in \Irr_{\bbK}(H)$ and $V \in \Irr_{\bbL}(H)$ with $X$ a submodule of $V|_{\bK[G]}$, or equivalently $V$ a submodule of $\bL \otimes_{\bbK} X$, then the three options from Theorem \ref{thm:facts} are summarised in Table \ref{tab:thewaysLHSsquare}. The three options from Proposition \ref{prop:algindressummary} recover these three cases: Case 1 is Case (B), Case 3(a) is Case (C) and Case 3(b) is Case (A).

\begin{table}
\caption{Induction and Restriction Possibilities Between $\bK[H]$ and $\bL[H]$}
\label{tab:thewaysLHSsquare}
\begin{center}
\def\arraystretch{1.3}
\begin{tabular}{ c c c c } 
    \toprule
Case & \textbf{A} & \textbf{B} & \textbf{C} \\
\midrule
 $\bL \otimes_{\bbK} X$ & $V$ & $V \oplus \overline{V}$ & $V \oplus V$ \\ 
 $V|_{\bK[H]}$ & $X \oplus X$ & $X$ & $X$\\ 
 $V \cong \overline{V}$? & Yes & No & Yes \\ 
  $m_c(V)$ & 1 & 1 & 2 \\
  $\dim_{\bbL} \End_{\bL[H]}(V) : \dim_{\bbK} \End_{\bK[H]}(X)$ & $1:1$ & $1:2$ & $1:4$ \\ 
  \bottomrule
\end{tabular}
\end{center} 
\end{table}

\subsection{The Extensions $\bK[H] \hookrightarrow \bK[G]$ and $\bL[H] \hookrightarrow \bL[G]$}\label{sect:indextwosubgroup}
Let us describe everything for $\bK[H] \hookrightarrow \bK[G]$, the case for $\bL[H] \hookrightarrow \bL[G]$ obtained by replacing $\bK$ by $\bL$ everywhere.

Write $\pi$ for the $1$-dimensional sign representation of $G$ over $H$, which takes value $1$ on $H$ and $-1$ on $G \setminus H$. The duality on $\Mod_{\bK[G]}$ is 
\[
(-) \otimes \pi \colon \Mod_{\bK[G]} \rightarrow \Mod_{\bK[G]}.
\]
It is direct to see that $\Res(V) = \Res(V \otimes \pi)$. The duality on $\Mod_{\bK[H]}$ is
\[
{}^\omega(-) \colon \Mod_{\bK[H]} \rightarrow \Mod_{\bK[H]},
\]
where ${}^{\omega} W$ is $W$ as a $\bK$-vector space and action of $h \in H$ by the action of $\omega^{-1} h \omega$ on $W$. There is an isomorphism of $\bK[G]$-modules
\[
    \bK[G] \otimes_{\bK[H]} {}^{\omega} W \xrightarrow{\sim} \bK[G] \otimes_{\bK[H]} W, \qquad g \otimes w \mapsto \omega^{-1} g \omega \otimes w
\]
and so $\Ind(W) = \Ind({}^{\omega} W)$. The following shows that the other properties also hold when $\charfield(\bK) \neq 2$.

\begin{prop}\label{prop:resindKHG}
Suppose that $W \in \Mod_{\bK[H]}$, $U \in \Mod_{\bK[G]}$, and $\charfield(K) \neq 2$. Then 
\begin{align*}
    \Res(\Ind(W)) &\cong W \oplus {}^{\omega} W, \\
    \Ind(\Res(U)) &\cong U \oplus (U \otimes \pi).
\end{align*}
\end{prop}

\begin{proof}
    The first isomorphism is given by the $K[H]$-module identifications
    \[  
        \bK[G] \otimes_{\bK[H]} W \equiv (\bK[H] \oplus \bK[H] \omega) \otimes_{\bK[H]} W \equiv W \oplus {}^{\omega} W.
    \]
    The second is given by the $K[G]$-module identification
    \[
        \bK[G] \otimes_{\bK[H]} U \xrightarrow{\sim} \bK[G/H] \otimes_{\bbK} U, \qquad g \otimes u \mapsto g \otimes g(u)
    \]
    and noting that the representation $\bK[G/H] = \mathbbm{1} \oplus \pi$ as $\bK[G]$-modules because $\charfield(\bK) \neq 2$.
\end{proof}

\begin{remark}
    Proposition \ref{prop:resindKHG} is false when $\charfield(\bK) \neq 2$: if it were true, then by Proposition \ref{prop:algindressummary} $\Ind(W)$ would always be semisimple for any simple  $W \in \Mod_{\bK[H]}$. However, in characteristic $2$ this isn't true when one takes $W$ to be the trivial representation.
\end{remark}

\subsection{The Extension $\bK[G] \hookrightarrow \bL_{\star \hspace{-0.1em}}[G]$} As in the previous section, we have the duality
\[
(-) \otimes \pi \colon \Mod_{\bK[G]} \rightarrow \Mod_{\bK[G]}.
\]
We have that the $\bK[G]$-module $\pi$ maps under $\Ind$ to the trivial representation of $\bL_{\star \hspace{-0.1em}}[G]$, by \cite[Thm.\ A]{TAYSchur} and the fact that $\Ind(\pi)$ restricts to the trivial representation as a $\bL[H]$-module. Therefore because $\Ind = \bL \otimes_{\bbK} -$ is a tensor functor (see \cite[\S4.1]{TAYSchur}), $\Ind(U) = \Ind(U \otimes \pi)$.

For $\Mod_{\bL_{\star \hspace{-0.1em}}[G]}$, we have a $\bK$-algebra automorphism $\overline{\cdot} \colon \bL_{\star \hspace{-0.1em}}[G] \rightarrow \bL_{\star \hspace{-0.1em}}[G]$ sending $\lambda g \mapsto \overline{\lambda} g$ which is well-defined because $\Gamma$ is abelian. We can therefore precompose any $\bL_{\star \hspace{-0.1em}}[G]$-module to with $\overline{\cdot}$, and to obtain the duality
\[
\overline{(-)} \colon \Mod_{\bL_{\star \hspace{-0.1em}}[G]} \rightarrow \Mod_{\bL_{\star \hspace{-0.1em}}[G]}.
\]
As the automorphism $\overline{\cdot}$ of $\bL_{\star \hspace{-0.1em}}[G]$ fixes $\bK[G]$, we have that $\Res(X) = \Res(\overline{X})$ for any $X \in \Mod_{\bL_{\star \hspace{-0.1em}}[G]}$. We can again check that the other properties hold when $\charfield(\bK) \neq 2$.

\begin{prop}\label{prop:resindKGLG}
Suppose that $U \in \Mod_{\bK[G]}$, $X \in \Mod_{\bL_{\star \hspace{-0.1em}}[G]}$, and $\charfield(\bK) \neq 2$. Then 
\begin{align*}
    \Res(\Ind(U)) &\cong U \oplus (U \otimes \pi), \\
    \Ind(\Res(X)) &\cong X \oplus \overline{X}.
\end{align*}
\end{prop}

\begin{proof}
    For the first isomorphism, $\bL \cong \bK[\Gamma]$ as a $\bK[G]$-module because $\bL/\bK$ is Galois, and $\bK[\Gamma] = \mathbbm{1} \oplus \pi$ because $\charfield(\bK) \neq 2$. Therefore
    \[
        \Res(\Ind(U)) = \bL \otimes_{\bbK} U = \bK[\Gamma] \otimes_{\bbK} U = U \oplus (U \otimes \pi).
    \]
    The second isomorphism is given by
    \[
        \bL \otimes_{\bbK} X \rightarrow X \oplus \overline{X}, \qquad \lambda \otimes x \mapsto (\lambda x, \overline{\lambda} x),
    \]
    which is easily verified to be $\bL_{\star \hspace{-0.1em}}[G]$-linear. To show it is bijective, by dimension considerations it suffices to show surjectivity. To this end, let $e = \sum_i \alpha_i \otimes \beta_i \in \bL \otimes_{\bbK} \bL$ be the unique element which maps to $(1,0)$ under the Galois isomorphism $\bL \otimes_{\bbK} \bL \rightarrow \bL \oplus \bL$, $a \otimes b \mapsto (ab,\overline{a} b)$. Then for any $x \in X$, 
    \[
        \sum_i \alpha_i \otimes \beta_i x \mapsto \left(\sum_i \alpha_i \beta_i x, \sum_i \overline{\alpha_i} \beta_i x \right) = (x,0).
    \]
    Similarly considering $(0,1)$ we can obtain any element $(0,x)$ and hence the map is surjective.
\end{proof}

\section{Analysis of the Diamond}\label{sect:diamondanalysis}

We continue with the notation from above, so $\bL / \bK$ is a Galois extension of degree two, $H \leq G$ is a Real structure on $H$, and $\omega$ is a fixed element of $G \setminus H$. Furthermore:

\vspace{0.5em}
\begin{center}
\emph{From now on in this chapter we additionally assume that $\charfield(\bK) \neq 2$.}
\end{center}
\vspace{0.0em}

\subsection{Group Actions on Sets of Irreducible Representations}

Let us start with the following version of the diamond~\eqref{eqn:diamond}, where we restrict to finite-dimensional irreducible representations:
\begin{equation}\label{eqn:irrdiamond}
\begin{tikzcd}
	& {  \Irr_{\bbK}(G) \, \rotatebox[origin=c]{90}{$\circlearrowleft$} \, C_2} & \\
  {{\Irr_{\bbL}(G)} \, \rotatebox[origin=c]{90}{$\circlearrowleft$} \, K_4}
  &
  {\Irr_{\bbK}(H) \, \rotatebox[origin=c]{90}{$\circlearrowleft$} \, C_2}
    &
    { {\Irr_{\bbL}^{\hspace{0.1em} \star}(G)}\, \rotatebox[origin=c]{90}{$\circlearrowleft$} \, C_2} \\
     & {\Irr_{\bbL}(H) \, \rotatebox[origin=c]{90}{$\circlearrowleft$} \, K_4}
	\arrow[tail reversed, from=1-2, to=2-1]
	\arrow[tail reversed, from=1-2, to=2-3]
	\arrow[tail reversed, from=2-2, to=1-2]
	\arrow[tail reversed, from=2-2, to=3-2]
	\arrow[tail reversed, from=3-2, to=2-1]
	\arrow[tail reversed, from=3-2, to=2-3]
\end{tikzcd}
\end{equation}

To describe the actions of the groups $K_4 = \{1,a,b,c\}$ and $C_2 = \{1,\epsilon\}$, we describe how these groups act on the sets of (isomorphism classes of) all representations, which restrict to the actions on the sets of (isomorphism classes of) irreducible representations above.

First, pick $Z = (Z,\rho) \in \Rep_{\bbL}(G)$. We can obtain three new representations:
\begin{itemize}
\item the Galois twist $\overline{Z}$ (as described in Section \ref{sect:chap3notation}), with the same action of $G$,
\item the twist ${}^{\omega}{Z} = (Z,{}^{\omega}\rho)$ where ${}^{\omega}\rho (x) = \rho (\omega^{-1} x \omega)$ for all $x\in G$. 
\item the tensor product $Z \otimes \pi \coloneqq Z \otimes_{\bbL} \pi$ with the sign character $\pi : G \to \{ \pm 1\}$. 
\end{itemize}  
The twist by $\omega$ is an inner twist, hence $\rho(\omega) \colon Z \xrightarrow{\sim} {}^{\omega}{Z}$ is an isomorphism and this operation gives nothing new. The remaining two operations commute, and we thus have an action of $K_4$:
\[
    a * Z = \overline{Z}, \qquad b*Z = Z \otimes \pi, \qquad c * Z = \overline{Z} \otimes \pi.
    \]

Consider now $V \in \Rep_{\bbL}(H)$. Here $\pi$ is not available, so the group $K_4$ acts by (cf.\ Section~\ref{sect:C2actionIrrH})
\[
    a * V = \overline{V}, \qquad b * V = {}^{\omega} V, \qquad c * V = \epsilon * V = {}^{\omega}\overline{V},
\]

As is true of any transitive $K_4$-set, the structure of the $K_4$-set $\{V, \overline{V}, {}^{\omega} V, {}^{\omega} \overline{V}\}$ is uniquely determined by $\Stab_{K_4}(V)$. Similarly, the $K_4$-set $\{Z, \overline{Z}, Z \otimes \pi, \overline{Z} \otimes \pi\}$ is determined by $\Stab_{K_4}(Z)$.

Now pick $U  \in \Rep_{\bbK}(G)$. Here the Galois twist is no longer available, while the twist by $\omega$ similarly produces an isomorphic representation. Hence, we have a $C_2$-action
\[
    \epsilon  * U =  U \otimes \pi \coloneqq U \otimes_{\bbK} \pi.
    \]

    Now consider $X  \in \Rep_{\bbK}(H)$. Here the Galois twist and the sign character are off limits.
All that remains is the twist by $\omega$, which yields a $C_2$-action
\[
    \epsilon  * X =  {}^{\omega}X.
    \]

    Finally, let us examine $W = (W,\rho) \in {\Irr_{\bbL}^{\hspace{0.1em} \star}(G)}$. All three operations appear available.
    However, we first note that $W \cong W \otimes_{\bbL} \pi$, because the semilinear sign representation $(\bbL, \pi)$ is isomorphic to the trivial semilinear representation $(\bbL, 1)$.
    Indeed, pick any non-zero ``imaginary'' element, that is, an $\alpha\in \bbL$ with $\overline{\alpha} = - \alpha \neq 0$. Then multiplication by $\alpha$ is an isomorphism $(\bbL, 1)\to (\bbL, \pi)$.


    What remains are the twist ${}^{\omega} W$ by $\omega$, and the conjugate representation $\overline{W}$, which defines a semilinear representation with the same action of $G$ because $\Gal(\bL / \bK)$ is abelian: for any $w \in W$, $\alpha \in \bL$, and $g \in G$,
    \[
    \rho(g)(\alpha * v) = \rho(g)(\overline{\alpha}v) = \sigma(g)(\overline{\alpha}) \rho(g)(v) = \overline{\sigma(g)(\alpha)} \rho(g)(v) = \sigma(g)(\alpha) * \rho(g)(v).
    \]
    
    Although both of ${}^{\omega} W$ and $\overline{W}$ differ from $W$ in general, they turn out to be the same: $\rho (\omega)$ can be viewed as an isomorphism $\overline{W} \to {}^{\omega}W$. Hence, we have only a $C_2$-action
\[
    \epsilon  * W =  {}^{\omega}W = \overline{W}.
    \]

\subsection{Propagation}\label{sect:propagation}

Let $V \in \Irr_{\bbL}(H)$. From this we can ``propagate out'' along the diamond~\eqref{eqn:irrdiamond} and define: \vspace{0.5em}
\begin{itemize}
    \item $W \in \Irr_{\bbL}^{\hspace{0.1em} \star}(G)$, an irreducible submodule of $\Ind_{\bL[H]}^{\bL_{\star \hspace{-0.1em}}[G]}(V)$, \vspace{0.5em}
    \item $X \in \Irr_{\bbK}(H)$, an irreducible submodule of $V|_{\bK[H]}$, \vspace{0.5em}
    \item $Z \in \Irr_{\bbL}(G)$, an irreducible submodule of $\Ind_{\bL[H]}^{\bL[G]}(V)$, \vspace{0.5em}
    \item $U \in \Irr_{\bbK}(G)$, an irreducible submodule of $Z|_{\bK[G]}$.
\end{itemize}

Using these modules, we can restrict the diamond~\eqref{eqn:diamond} to orbits under the corresponding groups.
The meaning of the labelling of the arrows will become clear in the next section.
\begin{equation}\label{eqn:labelleddiamond}
\begin{tikzcd}
	& {  \{U,U\otimes \pi \} \, \rotatebox[origin=c]{90}{$\circlearrowleft$} \, C_2} & \\
  {\{Z,\overline{Z},Z\otimes \pi ,\overline{Z}\otimes \pi  \} \, \rotatebox[origin=c]{90}{$\circlearrowleft$} \, K_4}
  &
  { \{X, {}^{\omega}X\} \, \rotatebox[origin=c]{90}{$\circlearrowleft$} \, C_2}
    &
    {\{W, \overline{W} \} \, \rotatebox[origin=c]{90}{$\circlearrowleft$} \, C_2} \\
    & {\{V,\overline{V},{}^{\omega}V,{}^{\omega}\overline{V} \} \, \rotatebox[origin=c]{90}{$\circlearrowleft$} \, K_4}
	\arrow["3"', tail reversed, from=1-2, to=2-1]
	\arrow["4", tail reversed, from=1-2, to=2-3]
	\arrow["4", tail reversed, from=2-2, to=1-2]
	\arrow["3"', tail reversed, from=2-2, to=3-2]
	\arrow["4", tail reversed, from=3-2, to=2-1]
	\arrow["3"', tail reversed, from=3-2, to=2-3]
\end{tikzcd}
\end{equation}

We have used four of the double arrows in the diamond~\eqref{eqn:irrdiamond} to define these modules. Equally, we could have used the 
remaining two, as clear from the next fact.
\begin{lemma}
    $U$ is a submodule of $\Ind_{\bK[H]}^{\bK[G]}(X)$ and $W$ is a submodule of $\Ind_{\bK[G]}^{\bL_{\star \hspace{-0.1em}}[G]}(U)$.
\end{lemma}

\begin{proof}
    For the first claim, first note that
    \[
    \left(Z, \Ind_{\bK[H]}^{\bL[G]}(X)\right) \neq 0,
    \]
    as $Z$ is a submodule of $\Ind_{\bL[H]}^{\bL[G]}(V)$ and $V \subset \Ind_{\bK[H]}^{\bL[H]}(X)$ by Frobenius reciprocity. Then 
    \[
        \left(Z, \Ind_{\bK[G]}^{\bL[G]}\left(\Ind_{\bK[H]}^{\bK[G]}(X)\right) \right) \neq 0
    \]
    so again by Frobenius reciprocity
    \[
    \left(Z|_{\bK[G]}, \Ind_{\bK[H]}^{\bK[G]}(X)\right) \neq 0.
    \]
    As $Z|_{\bK[G]}$ is a direct sum of copies of $U$ the first claim follows.

    For the second claim, we have that
    \[
    \left(W, \Ind_{\bK[H]}^{\bL_{\star \hspace{-0.1em}}[G]}(X)\right) \neq 0,
    \]
    as $V \subset \Ind_{\bK[H]}^{\bL[H]}(X)$ by Frobenius reciprocity and $W \subset \Ind_{\bL[H]}^{\bL_{\star \hspace{-0.1em}}[G]}(V)$. We have that
    \[
        \left(W, \Ind_{\bK[G]}^{\bL_{\star \hspace{-0.1em}}[G]}\left(\Ind_{\bK[H]}^{\bK[G]}(X)\right) \right) \neq 0,
    \]
    that $\Ind_{\bK[H]}^{\bK[G]}(X)$ is a direct sum of $U$ and $U \otimes \pi$ in equal multiplicity, and that
    \[
    \Ind_{\bK[G]}^{\bL_{\star \hspace{-0.1em}}[G]}(U) = \Ind_{\bK[G]}^{\bL_{\star \hspace{-0.1em}}[G]}(U \otimes \pi),
    \]
    hence
    \[
    \left(W, \Ind_{\bK[G]}^{\bL_{\star \hspace{-0.1em}}[G]}(U) \right) \neq 0,
    \]
    which establishes the claim.
\end{proof}

\subsection{The type and edge labels}
Each pair of modules adjacent in the diamond~\eqref{eqn:labelleddiamond} falls into one of the four cases 1, 2, 3(a) and 3(b) of Section \ref{sect:fourfoldway}. For three of these pairs there are only three possibilities: we label those instead by A, B and C as in Tables \ref{tab:thewaysRHSsquare} and \ref{tab:thewaysLHSsquare}. The double edges of the diamond~\eqref{eqn:labelleddiamond} are marked with the number of possibilities for each edge
and
we can record the induction and restriction behaviour of this family of modules as an element of
\[
T \coloneqq \{\text{A,\:B,\:C}\}^3 \times \{\text{1,\:2,\:3(a),\:3(b)}\}^3
\]
corresponding to the ordered extensions
\[
(\bK[G] \subset \bL[G], \: \bK[H] \subset \bL[H], \: \bL[H] \subset \bL_{\star \hspace{-0.1em}}[G]) \times (\bL[H] \subset \bL[G], \: \bK[H] \subset \bK[G], \: \bK[G] \subset \bL_{\star \hspace{-0.1em}}[G]).
\]
The ordering here is chosen to match the ordering obtained by reading the diamond from left to right.

\begin{defn}
The \emph{type} of $V$, $\sT(V)$, is the associated element of $T$.
\end{defn}

\begin{lemma}\label{lem:typeindependence}
    The type of $V$ is independent of the choices of the irreducible representations $X$, $Z$, $W$, $U$, and depends only on the $K_4$-orbit of $V$.
\end{lemma}

\begin{proof}
    From the explicit definition of each action, on any orbit of Diagram \eqref{eqn:labelleddiamond} there is an identification of endomorphism rings. Furthermore, for any extension of Diagram \eqref{eqn:labelleddiamond}, irreducible representations in the same orbit will have corresponding irreducible representations in the same orbit. Therefore, for any irreducible representation, and any extension of Diagram \eqref{eqn:labelleddiamond}, the endomorphism ring ratios of Proposition \ref{prop:algindressummary} will be constant on orbits, as will the properties of being fixed by the $C_2$-action of Proposition \ref{prop:algindressummary} (because $K_4$ and $C_2$ are abelian), and therefore the type for each extension will also be constant on orbits.    
\end{proof}

\begin{remark}
    We have chosen to start with $V$, but we may equally define the type of any other irreducible representation from any of the other four categories of Diagram \eqref{eqn:diamond}, which similarly will be independent of the choices involved and depend only on its orbit. We give an alternative point of view on the type in terms of the notion of an \emph{S-block} in Section \ref{sect:blockinterpretation} below.
\end{remark}

A priori, many possible types can occur, as the set $T$ has order $3^3 \cdot 4^3 = 1728$. The main result of this section is the following classification, which shows that only 24 possibilities can occur.

\begin{thm}\label{thm:24ways}
The possible types for $V$ that can occur are those 24 types listed in Table \ref{tab:24types}.
\end{thm}

In Section \ref{sect:realise24ways} we consider the question of realisation of the 24 types of Table \ref{tab:24types}. In the rest of this section we prove Theorem \ref{thm:24ways}. First we make some useful observations.

\begin{lemma}\label{lem:UnotUpi}
    If $U \not\cong U \otimes \pi$, then the type of $V$ is one of
    \emph{\begin{itemize}
        \item \textbf{(A, A, A, 1, 1, 2)},
        \item \textbf{(B, B, B, 1, 1, 2)},
        \item \textbf{(C, C, C, 1, 1, 2)}.
    \end{itemize}}
\end{lemma}
\begin{proof}
    If $U \not\cong U \otimes \pi$, then $\bK[H] \subset \bK[G]$ is in Case 1, and $\bK[G] \subset \bL_{\star \hspace{-0.1em}}[G]$ is in Case 2. As $\Res(Z \otimes \pi) = \Res(Z) \otimes \pi$, we have that $Z \not\cong Z \otimes \pi$, so $\bL[H] \subset \bL[G]$ is in Case 1. From this, we have that $(W,W) = (U,U) = (X,X)$ and $(Z,Z) = (V,V)$. The three remaining cases for each of the remaining extensions are determined by the ratio $(V,V) : (X,X)$ on endomorphism rings, and a choice for one extension determines the other two. The three possibilities are listed above.
\end{proof}

\begin{lemma}\label{lem:WnotomegaW}
    If $X \not\cong {}^{\omega}X$, then the type of $V$ is one of
    \emph{\begin{itemize}
        \item \textbf{(A, A, B, 2, 2, 3(b))},
        \item \textbf{(B, B, B, 2, 2, 1)},
        \item \textbf{(C, C, B, 2, 2, 3(a))}.
    \end{itemize}}
\end{lemma}
\begin{proof}
    If $X \not\cong {}^{\omega}X$, then $\bK[H] \subset \bK[G]$ is in Case 2, and $(U,U) = (X,X)$. As $X \not\cong {}^{\omega}X$, $\{V, \overline{V}\} \neq \{{}^{\omega}V, {}^{\omega}\overline{V}\}$, so $V \not\cong {}^{\omega} \overline{V}$, meaning that $\bL[H] \subset \bL_{\star \hspace{-0.1em}}[G]$ is in Case B and $V \not\cong {}^{\omega} \overline{V}$, meaning that $\bL[H] \subset \bL[G]$ is in Case 2. So the type is (?, ?, B, 2, 2, ?), and $(W,W) = (V,V) = (Z,Z)$. We now consider the possible options for $\bK[H] \subset \bK[G]$, and show these determine the rest.
    \begin{itemize}
        \item In Case A, $2(X,X) = (V,V)$, so $2(U,U) = (W,W)$, and $\bK[G] \subset \bL_{\star \hspace{-0.1em}}[G]$ is in Case 3(b). Furthermore, $2(U,U) = (Z,Z)$, so $\bK[G] \subset \bL[G]$ is in Case A.
        \item In Case B, $(X,X) = (V,V)$, so $(U,U) = (W,W)$, and as $U \cong U \otimes \pi$, $\bK[G] \subset \bL_{\star \hspace{-0.1em}}[G]$ is in Case 1. Furthermore, $(U,U) = (Z,Z)$, so $\bK[G] \subset \bL[G]$ is in Case B.
        \item In Case C, $(X,X) = 2(V,V)$, so $(U,U) = 2(W,W)$, and $\bK[G] \subset \bL_{\star \hspace{-0.1em}}[G]$ is in Case 3(a). Furthermore, $(U,U) = 2(Z,Z)$, so $\bK[G] \subset \bL[G]$ is in Case C. \qedhere
    \end{itemize}
\end{proof}

\subsection{Proof of Theorem \ref{thm:24ways}} We now being our analysis of the possible types of $V$ in earnest. It is essentially a case analysis, indexed first by the extension $\bK[H] \subset \bL[H]$ and then by $\bK[H] \subset \bK[G]$.

\vspace{0.5em}
\textcolor{gray}{\hrule}

\vspace{0.9em}
\begin{center}
    \large{\textbf{Case A} (?, A, ?, ?, ?, ?)}
\end{center}
\[
 V \cong \overline{V}, \ \Ind(X) = V, \ \Res(V) = X \oplus X, \ (X,X) = \dim_{\bbL} \End_{\bL[H]}(V).
\]
\vspace{-0.3em}
\textcolor{gray}{\hrule}
\vspace{0.5em}

\vspace{0.5em}
\noindent \textbf{Case A.1} (?, A, ?, ?, 1, ?)

From Lemma \ref{lem:UnotUpi}, the type of $V$ must be \textbf{(A, A, A, 1, 1, 2)}.

\vspace{0.5em}
\noindent \textbf{Case A.2} (?, A, ?, ?, 2, ?)

From Lemma \ref{lem:WnotomegaW}, the type of $V$ must be \textbf{(A, A, B, 2, 2, 3(b))}.

\vspace{0.5em}
\noindent \textbf{Case A.3(a)} (?, A, ?, ?, 3(a), ?)

In this case we have that
\[
U \cong U \otimes \pi, \ X \cong {}^{\omega} X,\ \Ind(X) = U \oplus U, \ \Res(U) = X, \ (X,X) = 2(U,U).
\]
Because $X \cong {}^{\omega} X$ and $\Ind(X) \cong V$, $V \cong {}^{\omega} V$, and therefore, as $V \cong \overline{V}$, $\Stab_{K_4}(V) = K_4$. In particular, $\bL[H] \subset \bL_{\star \hspace{-0.1em}}[G]$ is either in Case A or Case C. In Case C, we would have $(W,W) = 4 \dim_{\bbL} \End_{\bL[H]}(V)$, and thus would have $(W,W) = 4 (X,X) = 8(U,U)$, which cannot occur. Therefore, we must be in Case A, with type (?, A, A, ?, 3(a), ?) and $(W,W) = \dim_{\bbL} \End_{\bL[H]}(V)$. In particular,
\[
(W,W) = \dim_{\bbL} \End_{\bL[H]}(V) = (X,X) = 2(U,U),
\]
and so $\bK[G] \subset \bL_{\star \hspace{-0.1em}}[G]$ is in Case 3(b) and the type is (?, A, A, ?, 3(a), 3(b)).

There are then three options for the extension $\bK[G] \subset \bL[G]$.
\begin{itemize}
\item In Case A, type (A, A, A, ?, 3(a), 3(b)), $(U,U) = \dim_{\bbL}\End_{\bL[G]}(Z)$. Therefore
    \[
    \dim_{\bbL}\End_{\bL[H]}(V) = (X,X) = 2(U,U) = 2 \dim_{\bbL}\End_{\bL[G]}(Z),
    \]
    hence $\bL[H] \subset \bL[G]$ is in Case 3(a). The type is \textbf{(A, A, A, 3(a), 3(a), 3(b))}.
\item In Case B, $(U,U) = 2\dim_{\bbL}\End_{\bL[G]}(Z)$, hence
\[
\dim_{\bbL}\End_{\bL[H]}(V) = (X,X) = 2(U,U) = 4 \dim_{\bbL}\End_{\bL[G]}(Z),
\]
which cannot occur.
\item In Case C, $(U,U) = 4\dim_{\bbL}\End_{\bL[G]}(Z)$, hence
\[
\dim_{\bbL}\End_{\bL[H]}(V) = (X,X) = 2(U,U) = 8 \dim_{\bbL}\End_{\bL[G]}(Z),
\]
which cannot occur.
\end{itemize}

\vspace{0.5em}
\noindent \textbf{Case A.3(b)} (?, A, ?, ?, 3(b), ?)

In this case we have that 
\[
U \cong U \otimes \pi, \ X \cong {}^{\omega} X, \ \Ind(X) = U, \ \Res(U) = X \oplus X, \ 2(X,X) = (U,U).
\]
Because $X \cong {}^{\omega} X$ and $\Ind(X) \cong V$, $V \cong {}^{\omega} V$, and therefore, as $V \cong \overline{V}$, $\Stab_{K_4}(V) = K_4$. In particular, $\bL[H] \subset \bL_{\star \hspace{-0.1em}}[G]$ is either in Case A or Case C. Let us first observe that being A or C determines the case for $\bK[G] \subset \bL_{\star \hspace{-0.1em}}[G]$: 

\begin{itemize}
    \item In Case A, $(W,W) = \dim_{\bbL} \End_{\bL[H]}(V)$, so $(U,U) = 2(X,X) = 2 (W,W)$ and $\bK[G] \subset \bL_{\star \hspace{-0.1em}}[G]$ is in Case 3(a). The type is (?, A, A, ?, 3(b), 3(a)).
    \item In Case C, $(W,W) = 4 \dim_{\bbL} \End_{\bL[H]}(V)$, and thus $(W,W) = 4 (X,X) = 2(U,U)$, so $\bK[G] \subset \bL_{\star \hspace{-0.1em}}[G]$ is in Case 3(b). The type is (?, A, C, ?, 3(b), 3(b)).
\end{itemize}

Note also that in either case $(V,V) = 2 \dim_{\bbL} \End_{\bL[H]}(V) = 2(X,X) = (U,U)$. We can use this to see that the three cases for $\bK[G] \subset \bL[G]$ determine those for $\bL[H] \subset \bL[G]$: Case 2 cannot occur for $\bL[H] \subset \bL[G]$, as $V \cong {}^{\omega}V$, and the three remaining cases are determined by $(V,V):(Z,Z)$, which is determined by the case for $\bK[G] \subset \bL[G]$. The three options are (B, A, ?, 1, 3(b), ?), (C, A, ?, 3(a), 3(b), ?), and (A, A, ?, 3(b), 3(b), ?).

Pairing these two options for the extension $\bL[H] \subset \bL_{\star \hspace{-0.1em}}[G]$ with the three options for $\bK[G] \subset \bL[G]$, we obtain six options:
\begin{itemize}
     \item \textbf{(B, A, A, 1, 3(b), 3(a))},
     \item \textbf{(B, A, C, 1, 3(b), 3(b))},
     \item \textbf{(C, A, A, 3(a), 3(b), 3(a))},
     \item \textbf{(C, A, C, 3(a), 3(b), 3(b))},
     \item \textbf{(A, A, A, 3(b), 3(b), 3(a))},
     \item \textbf{(A, A, C, 3(b), 3(b), 3(b))},
\end{itemize}

\vspace{0.5em}
\textcolor{gray}{\hrule}

\vspace{0.9em}
\begin{center}
    \large{\textbf{Case B} (?, B, ?, ?, ?, ?)}
\end{center}
\[
 V \not\cong \overline{V}, \ \Ind(X) = V \oplus \overline{V}, \ \Res(V) = X, \ (X,X) = 2\dim_{\bbL} \End_{\bL[H]}(V).
\]
\vspace{-0.3em}
\textcolor{gray}{\hrule}
\vspace{0.5em}

\vspace{0.5em}
\noindent \textbf{Case B.1} (?, B, ?, ?, 1, ?)

From Lemma \ref{lem:UnotUpi}, the type of $V$ must be \textbf{(B, B, B, 1, 1, 2)}.

\vspace{0.5em}
\noindent \textbf{Case B.2} (?, B, ?, ?, 2, ?)

From Lemma \ref{lem:WnotomegaW}, the type of $V$ must be \textbf{(B, B, B, 2, 2, 1)}.

\vspace{0.5em}
\noindent \textbf{Case B.3(a)} (?, B, ?, ?, 3(a), ?)

In this case we have that 
\[
U \cong U \otimes \pi, \ X \cong {}^{\omega} X, \ \Ind(X) = U \oplus U, \ \Res(U) = X, \ (X,X) = 2(U,U).
\]
We cannot have that $\bL[H] \subset \bL_{\star \hspace{-0.1em}}[G]$ is in Case C, otherwise $4(U,U) = 2(X,X) = 2(V,V) = (W,W)$, which cannot happen for the extension $\bK[G] \subset \bL_{\star \hspace{-0.1em}}[G]$. Therefore we have two options for $\bL[H] \subset \bL_{\star \hspace{-0.1em}}[G]$:
\begin{itemize}
    \item In Case A, $(V,V) = 2(W,W)$, so as $(V,V) = 2(U,U)$, $(W,W) = (U,U)$, and $\bK[G] \subset \bL_{\star \hspace{-0.1em}}[G]$ is in Case 1, with type (?, B, A, ?, 3(a), 1). As $V \not\cong \overline{V}$ and $V \cong {}^{\omega} \overline{V}$, $\Stab_{K_4}(V) = \langle c \rangle$ and $V \not\cong {}^{\omega} V$ and $\bL[H] \subset \bL[G]$ is in Case 2, with $(Z,Z) = (V,V)$. Then $(Z,Z) = 2(U,U)$, so $\bK[G]\subset \bL[G]$ is in Case A, and the type is \textbf{(A, B, A, 2, 3(a), 1)}.
    \item In Case B, $(V,V) = (W,W)$ so as $(V,V) = 2(U,U)$, $(W,W) = 2(U,U)$, and $\bK[G] \subset \bL_{\star \hspace{-0.1em}}[G]$ is in Case 3(b), with type (?, B, B, ?, 3(a), 3(b)). If $Z \cong \overline{Z}$, then $\{V, {}^{\omega} V \} = \{\overline{V}, {}^{\omega}\overline{V}\}$, which is a contradiction, as $V \not\cong \overline{V}$ and $V \not\cong {}^{\omega}\overline{V}$. Therefore $Z \not\cong \overline{Z}$, and $\bK[G] \subset \bL[G]$ is in Case B. This forces $(V,V) = 2(Z,Z)$, so $\bL[H] \subset \bL[G]$ is in Case 3(a). The type is \textbf{(B, B, B, 3(a), 3(a), 3(b))}.
\end{itemize}

\vspace{0.5em}
\noindent \textbf{Case B.3(b)} (?, B, ?, ?, 3(b), ?)

In this we have that
\[
U \cong U \otimes \pi, \ X \cong {}^{\omega} X, \ \Ind(X) = U, \ \Res(U) = X \oplus X, \ 2(X,X) = (U,U).
\]
Suppose first that $W \cong \overline{W}$. Then $\{V, {}^{\omega} \overline{V} \} = \{\overline{V}, {}^{\omega}V\}$, so, as $V \not\cong \overline{V}$, $V \cong {}^{\omega} V$, and thus $\Stab_{K_4}(V) = \langle b \rangle$, so $V \not\cong {}^{\omega} \overline{V}$ and $\bL[H] \subset \bL_{\star \hspace{-0.1em}}[G]$ is in Case B. As $(U,U) = 2(W,W)$, $\bK[G] \subset \bL_{\star \hspace{-0.1em}}[G]$ is in Case 3(a), with type (?, B, B, ?, 3(b), 3(a)). If $Z \cong \overline{Z}$, then $\{V, {}^{\omega} V \} = \{\overline{V}, {}^{\omega}\overline{V}\}$, which is a contradiction, as $V \not\cong \overline{V}$ and $V \not\cong {}^{\omega}\overline{V}$. Therefore $Z \not\cong \overline{Z}$, and $\bK[G] \subset \bL[G]$ is in Case B. Therefore $(Z,Z) = 2(V,V)$, and $\bL[H] \subset \bL[G]$ is in Case 3(b). The type is \textbf{(B, B, B, 3(b), 3(b), 3(a))}.

If instead $W \not\cong \overline{W}$, then $\bK[G] \subset \bL_{\star \hspace{-0.1em}}[G]$ is in Case 1, so $(U,U) = (W,W)$. This means that $(W,W) = 2(V,V)$, and $\bL[H] \subset \bL_{\star \hspace{-0.1em}}[G]$ is in Case C. The type is (?, B, C, ?, 3(b), 1). From this, $\Stab_{K_4}(V) = \langle c \rangle$, hence $V \not\cong {}^{\omega} V$, and $\bL[H] \subset \bL[G]$ is in Case 2, and $2(Z,Z) = (U,U)$. Therefore $\bK[G] \subset \bL[G]$ is in Case C, and the type is \textbf{(C, B, C, 2, 3(b), 1)}.

\vspace{0.5em}
\textcolor{gray}{\hrule}

\vspace{0.9em}
\begin{center}
    \large{\textbf{Case C} (?, C, ?, ?, ?, ?)}
\end{center}
\[
 V \cong \overline{V}, \ \Ind(X) = V \oplus V, \ \Res(V) = X, \ (X,X) = 4\dim_{\bbL} \End_{\bL[H]}(V).
\]
\vspace{-0.3em}
\textcolor{gray}{\hrule}
\vspace{0.5em}

\vspace{0.5em}
\noindent \textbf{Case C.1} (?, C, ?, ?, 1, ?)

From Lemma \ref{lem:UnotUpi}, the type of $V$ must be \textbf{(C, C, C, 1, 1, 2)}.

\vspace{0.5em}
\noindent \textbf{Case C.2} (?, C, ?, ?, 2, ?)

From Lemma \ref{lem:WnotomegaW}, the type of $V$ must be \textbf{(C, C, B, 2, 2, 3(a))}.

\vspace{0.5em}
\noindent \textbf{Case C.3(a)} (?, C, ?, ?, 3(a), ?)

In this case we have that
\[
U \cong U \otimes \pi, \ X \cong {}^{\omega} X, \ \Ind(X) = U \oplus U, \ \Res(U) = X, \ (X,X) = 2(U,U).
\]
Because $X \cong {}^{\omega} X$, $V \cong {}^{\omega} V$, and therefore, as $V \cong \overline{V}$, $\Stab_{K_4}(V) = K_4$. In particular, $\bL[H] \subset \bL_{\star \hspace{-0.1em}}[G]$ is either in Case A or Case C. Let us first observe that this determines the case for $\bK[G] \subset \bL_{\star \hspace{-0.1em}}[G]$: 

$\bullet$ In Case A, $2(W,W) = (V,V)$, so as $(U,U) = (V,V)$, $(U,U) = 2 (W,W)$ and $\bK[G] \subset \bL_{\star \hspace{-0.1em}}[G]$ is in Case 3(a). The type is (?, C, A, ?, 3(a), 3(a)).

$\bullet$ In Case C, $(W,W) = 2 (V,V)$, so as $(U,U) = (V,V)$, $2(U,U) = (W,W)$ and $\bK[G] \subset \bL_{\star \hspace{-0.1em}}[G]$ is in Case 3(b). The type is (?, C, C, ?, 3(a), 3(b)).

As $(V,V) = (U,U)$, we see also that the three cases for $\bK[G] \subset \bL[G]$ determine those for $\bL[H] \subset \bL[G]$: Case 2 cannot occur for $\bL[H] \subset \bL[G]$, as $V \cong {}^{\omega}V$, and the three remaining cases are determined by $(V,V):(Z,Z)$, which is determined by the case for $\bK[G] \subset \bL[G]$. The three options are (B, C, ?, 1, 3(a), ?), (C, C, ?, 3(a), 3(a), ?), and (A, C, ?, 3(b), 3(a), ?).

Pairing these two options for the extension $\bL[H] \subset \bL_{\star \hspace{-0.1em}}[G]$ with the three options for $\bK[G] \subset \bL[G]$, we obtain six options:
\begin{itemize}
     \item \textbf{(B, C, A, 1, 3(a), 3(a))},
     \item \textbf{(B, C, C, 1, 3(a), 3(b))},
     \item \textbf{(C, C, A, 3(a), 3(a), 3(a))},
     \item \textbf{(C, C, C, 3(a), 3(a), 3(b))},
     \item \textbf{(A, C, A, 3(b), 3(a), 3(a))},
     \item \textbf{(A, C, C, 3(b), 3(a), 3(b))}.
\end{itemize}

\vspace{0.5em}
\noindent \textbf{Case C.3(b)} (?, C, ?, ?, 3(b), ?)
In this case we have that
\[
U \cong U \otimes \pi, \ X \cong {}^{\omega} X, \ \Ind(X) = U, \ \Res(U) = X \oplus X, \ 2(X,X) = (U,U).
\]
As the ratio $(U,U) : (X,X) : (V,V)$ along the central axis of the diamond is is $4:2:1$, and no pair has a ratio greater than $2$, we must have that $(Z,Z) = 2(V,V)$, $(U,U) = 2(Z,Z)$, $(W,W) = 2 (V,V)$, and $(U,U) = 2(W,W)$. Therefore the type is \textbf{(C, C, C, 3(b), 3(b), 3(a))}. \qedsymbol

\section{Behaviour of Each Type}

In Section \ref{sect:realise24ways} below, we consider which of the 24 types of Table~\ref{tab:24types} actually occur. Before that, in this section we record the behaviour of each type in Tables \ref{tab:24types2}, \ref{tab:24types4}, \ref{tab:24types5}, \ref{tab:24types5.5} and \ref{tab:24types6}. We will use this to identify the types in the examples of Section \ref{sect:realise24ways} below.

\subsection{Justification of Tables \ref{tab:24types2}, \ref{tab:24types4}, \ref{tab:24types5}, \ref{tab:24types5.5} and \ref{tab:24types6}}

Table \ref{tab:24types2}, Table \ref{tab:24types4}, and Table \ref{tab:24types6} each follow directly from Theorem \ref{thm:threeoptions}, Proposition \ref{prop:algindressummary} and Table \ref{tab:thewaysLHSsquare}. 

These also give the dimension ratios appearing in Table \ref{tab:24types5.5}. When all entries of a row of Table \ref{tab:24types5.5} are divisible by two (resp. four), this extra divisibility property of the dimensions follows from the fact that the dimensions of the endomorphism rings of Table \ref{tab:24types6} force the dimension of the representation of minimal dimension (which is always $V$) of a row to be a multiple of two (resp. four). The rows with minimal dimension equal to two are (2), (6), (7), (13), (16), (17), (18), (19), (22), (23), (24) and the rows with minimal dimension equal to four are (20) and (21).

Finally, Table \ref{tab:24types5} also follows from Theorem \ref{thm:threeoptions}, Proposition \ref{prop:algindressummary} and Table \ref{tab:thewaysLHSsquare}, the only part requiring further justification being the computation of $\Stab_{K_4}(Z)$ and $\Stab_{K_4}(V)$.

\subsubsection{Computation of $\Stab_{K_4}(Z)$}

To compute $\Stab_{K_4}(Z)$, note that if at least one of $a * Z \cong Z$ or $b * Z \cong Z$ is true, then $\Stab_{K_4}(Z)$ is determined by $a * Z$ and $b * Z$, which can be deduced from the type. The remaining cases to consider are (4), (5), (10), (18), (19), where $a * Z \not\cong Z$ or $b * Z \not\cong Z$, and so $\Stab_{K_4}(Z)$ could be either $1$ or $\langle c \rangle$. This can be deduced from whether or not $V \cong \overline{V}$, as follows.

If $V \cong \overline{V}$, then
\[
\Ind_{\bL[H]}^{\bL[G]}(V) \cong \Ind_{\bL[H]}^{\bL[G]}(\overline{V}),
\]
hence
\[
Z \oplus (Z \otimes \pi) \cong \overline{Z} \oplus (\overline{Z} \otimes \pi),
\]
so, as $Z \not\cong \overline{Z}$, $Z \cong \overline{Z} \otimes \pi$ and the stabiliser is $\langle c \rangle$.

Conversely, if $\overline{Z} \cong Z \otimes \pi$, then
\[
\Res_{\bL[H]}^{\bL[G]}(\overline{Z}) \cong \Res_{\bL[H]}^{\bL[G]}(Z \otimes \pi) = \Res(Z),
\]
so $V \cong \overline{V}$. 

Therefore in cases (4), (5), (18) and (19) the stabiliser is $\langle c \rangle$, and $\Stab_{K_4}(Z) = 1$ in case (10).

\subsubsection{Computation of $\Stab_{K_4}(V)$}

Whether or not $\overline{V} = a * V \cong V$, ${}^{\omega} V = b * V \cong V$, or $\epsilon * V = c * V \cong V$ can be read off from the type of $\bK[G] \subset \bL[G]$, $\bL[H] \subset \bL[G]$ and $\bL[H] \subset \bL_{\star \hspace{-0.1em}}[G]$ respectively.

\subsection{When $\bL$ Splits $H$ and $G$}\label{sect:LsplitsHG}

Certain types of the 24 of Theorem \ref{thm:24ways} appearing in Table \ref{tab:24types} can be ruled out depending on whether certain modules are absolutely irreducible. Indeed, from Table \ref{tab:24types6} we can see that
\begin{itemize}
    \item If $\End_{\bK[H]}(X) = \bK$, then there cannot be 3(a) in the 5th column,
    \item If $\End_{\bK[G]}(U) = \bK$, then there cannot be 3(b) in the 5th column, or 3(a) in the 6th column,
    \item If $\End_{\bL[H]}(V) = \bL$, then there cannot be 3(a) in the 4th column.
    \item If $\End_{\bL[G]}(Z) = \bL$, then there cannot be 3(b) in the 4th column.
\end{itemize}

For example, the number of possible types is at most
\begin{itemize}
    \item 12, when $\bK$ splits $H$,
    \item 10, when $\bK$ splits $G$,
    \item 5, when $\bK$ splits $H$ and $G$,
    \item 18, when $\bL$ splits $H$,
    \item 18, when $\bL$ splits $G$,
    \item 12, when $\bL$ splits $H$ and $G$.
\end{itemize}

Of particular interest is the final case when $\bL$ is a splitting field for $H$ and $G$. As described above, in this case 12 of the 24 types of Theorem \ref{thm:24ways} are ruled out, and we choose to give a alternative simpler name the remaining types.

\begin{defn}\label{def:twelvetypes}
For $V \in \Irr_{\bbL}(H)$, with $\End_{\bbL[H]}(V) = \bL$ and $\End_{\bbL[G]}(Z) = \bL$, the D-type of $V$, $\fD(V)$, is the roman numeral associated to the type as below. \vspace{0.5em}

\begin{minipage}[t]{0.48\textwidth}  % left column
\begin{tabular}{c l c}
I & \textbf{(A, A, A, 1, 1, 2)} & (1) \\
II & \textbf{(B, A, C, 1, 3(b), 3(b))} & (5) \\
II.V & \textbf{(B, A, A, 1, 3(b), 3(a))} & (4) \\
III & \textbf{(A, A, B, 2, 2, 3(b))} & (2) \\
IV & \textbf{(B, B, B, 1, 1, 2)} & (10) \\
V & \textbf{(A, B, A, 2, 3(a), 1)} & (12) \\
\end{tabular}
\end{minipage}%
\hfill
\begin{minipage}[t]{0.48\textwidth}  % right column
\begin{tabular}{c l c}
VI & \textbf{(C, B, C, 2, 3(b), 1)} & (15) \\
VII & \textbf{(B, B, B, 2, 2, 1)} & (11) \\
VIII & \textbf{(C, C, C, 1, 1, 2)} & (16) \\
IX & \textbf{(B, C, A, 1, 3(a), 3(a))} & (18)\\
IX.V & \textbf{(B, C, C, 1, 3(a), 3(b))} & (19)\\
X & \textbf{(C, C, B, 2, 2, 3(a))} & (17) \\
\end{tabular}
\end{minipage}
\end{defn}

We will see in Section \ref{sect:recoverDyson10} of Chapter \ref{chapter:tenfoldway} how these are consistent with the labelling of Dyson's tenfold way, with the two new cases that occur here labelled by II.V and IX.V. The motivation behind the naming of the new cases II.V and IX.V is that these D-types have the same L-type (see Section \ref{sect:forgetsemilineardata} below) as II and IX respectively. More precisely, under the map $\Psi$ of Section \ref{sect:forgetsemilineardata} below, $\Psi(\text{II}) = \Psi(\text{II.V})$ and $\Psi(\text{IX}) = \Psi(\text{IX.V})$. Note however that $\Psi$ \emph{is} injective when restricted to the other eight types.

There are other nice facts that happen when $\bL$ splits $H$ and $G$, such as a more explicit description of the endomorphism rings involved in cases A, B and C.

\begin{prop}\label{prop:endringwhensplit}
    Suppose that $V, W$ are as in Theorem \ref{thm:threeoptions}, and $\End_{\bL[H]}(V) = \bL$. Then the S-type of $V$ is characterised by the following conditions:
    \begin{enumerate}
        \item[\emph{(A)}] $\End_{\bL_{\star \hspace{-0.1em}}[G]}(W) = \bK$,
        \item[\emph{(B)}] $\End_{\bL_{\star \hspace{-0.1em}}[G]}(W) = \bL$,
        \item[\emph{(C)}] $\End_{\bL_{\star \hspace{-0.1em}}[G]}(W) = D$, where $D$ is a central division algebra over $\bK$ of degree $2$.
    \end{enumerate}
\end{prop}

\begin{proof}
    This follows directly from \cite[Prop.\ 8.28]{TAYSchur}.
\end{proof}

\begin{cor}
    When $\bL$ splits $H$ and $G$, the endomorphism rings for each type are described in Table \ref{tab:12types}.
\end{cor}

\begin{proof}
We know the first two columns must be $\bL$ from the fact that $\bL$ splits $H$ and $G$. We therefore know the dimensions of the remaining entries of each row by Table \ref{tab:24types6}. If the dimension is $2$, then the endomorphism ring is $\bL$ by Proposition \ref{prop:endringwhensplit}. Otherwise, from Proposition \ref{prop:endringwhensplit} the endomorphism is a central division algebra $D$ of dimension $4$ over $\bK$. A priori, these endomorphism rings are different, however in all but D-type IX.V the division algebra is the same.

Indeed, for any pair of adjacent columns among the final three, at least one of induction or restriction between the corresponding algebras preserves irreducibility, by Table \ref{tab:indres}, and therefore by functoriality there is a map from one endomorphism ring to the other. This has to be injective, because each endomorphism ring is a central simple algebra, and therefore an isomorphism by comparing the dimension.
\end{proof}

\begin{remark}
    The $\bK$-algebras $D_1$ and $D_2$ in D-type IX.V of Table \ref{tab:12types} are non-isomorphic in general. See Remark \ref{rem:differentdivalg} of Section \ref{sect:realiseIXV}.
\end{remark}

\section{Forgetting the Semilinear Data}\label{sect:forgetsemilineardata}

Given $V \in \Irr_{\bbL}(H)$, we can assign to $V$ the type of $V$, which we may consider as a map
\[
    \fT \colon  \Irr_{\bbL}(H) \rightarrow T,
\]
which by Theorem \ref{thm:24ways} lands in the subset $T_0 \subset T$ of the 24 possible types. Setting
\[
S \coloneqq \{\text{A,\:B,\:C}\}^2 \times \{\text{1,\:2,\:3(a),\:3(b)}\}^2,
\]
we also have a forgetful map
\[
\Psi \colon T \rightarrow S
\]
defined by forgetting the 3rd and 6th entry of an element of $T$. By construction this map remembers exactly the linear representation theoretic information about the type of $V$, which forms the left-hand side of the diamond (\ref{eqn:labelleddiamond}).

One of our original motivations for considering the type was to understand if the semilinear Schur index can always be determined the linear part of the type: the part obtainable without reference to semilinear representations. The following shows that this is not possible in general.

\begin{cor}
The map $\Psi|_{\fT(\Irr_{\bbL}(H))} \colon \fT(\Irr_{\bbL}(H)) \rightarrow S$ is not injective in general.
\end{cor}

In other words: it is impossible in general to recover the type of an irreducible $V \in \Irr_{\bbL}(H)$ from linear representation-theoretic data alone.

\begin{proof}
    Let $\bL / \bK$ be the extension $\bQ(i) / \bQ$, and let $G = \text{SmallGroup}(32,9)$, which has a unique subgroup $H$ isomorphic to $D_8 \times C_2$. Considering $G$ with the presentation
    \[
G := \langle a,b,c \mid 
a^4 = b^2 = c^4 = 1,\;
bab = c a c^{-1} = a^{-1},\;
c b c^{-1} = ab
\rangle
\]
the subgroup $H$ is given by $H = \langle a, b, c^2 \rangle$. The types of the two $2$-dimensional irreducible representations of $H$ over $\bQ(i)$ are (8) and (9), which have the same linear type. To see this, one can note that the field $\bQ(i)$ is a splitting field for $H$, as $H$ has exponent $4$, and therefore one can use Theorem \ref{thm:15waysforcyclotomic} to determine the type by computing Frobenius-Schur indicators. Alternatively, one can use the algorithm to compute the linear type found at \cite{GITHUB}.
\end{proof}

\subsection{The Linear Type}\label{sect:lineartype}

We include a repository \cite{GITHUB} containing magma code for computing the \emph{linear type} of $V \in \Irr_{\bbL}(H)$, or in other words, an algorithm for the function
\[
   \fL \coloneqq \colon \Psi \circ \fT \colon \Irr_{\bbL}(H) \rightarrow T_0 \rightarrow S.
\]
We note that the fibers of the map
\[
    \Psi|_{T_0} \colon T_0 \rightarrow S
\]
which have size at least two are:
\[
    \{(4), (5)\}, \quad \{(6), (7)\}, \quad \{(8), (9)\}, \quad \{(18), (19)\}, \quad \{(20), (21)\}, \quad \{(22), (23)\}.
\]
In particular, we have the following. 
\begin{cor}\label{cor:canrecoverSLtype}
Let $V \in \Irr_{\bbL}(H)$. If the linear type $\fL(V)$ satisfies
\[
\fL(V) \in \{(1), (2), (3), (10), (11), (12), (13), (14), (15), (16), (17), (24)\},
\]
then the type $\fT(V)$ can be uniquely recovered from $\fL(V)$.
\end{cor}

\section{Realisation of The 24 Types}\label{sect:realise24ways}

In this section we consider realisability of the 24 types of Theorem \ref{thm:24ways} and Table \ref{tab:24types}. We give explicit realisations of all but one type, leaving the realisation of type (21) as an open question. We give all our realisations over number fields, and in the twelve cases of Section \ref{sect:LsplitsHG} we further take $V \in \Irr_{\bbL}(H)$ to be such that $\End_{\bL[H]}(V) = \bL$ and $\End_{\bL[G]}(Z) = \bL$.

In light of Section \ref{sect:lineartype}, to determine the type of a representation, we typically first compute the linear type, which can either be done by hand or using the algorithm of Section \ref{sect:lineartype}. This either determines the type, or reduces the type to two possibilities.

These two possibilities are distinguished by the S-type of $V$, which to determine one has to typically either give an explicit extension of $V$ to a semilinear representation of $G$, or show by hand that no such extension is possible.

\subsection{Realisation of Type (1) / I}

Let $H \leq G$ be any Real structure, and let $\bL / \bK$ be any degree two extension with $\charfield(\bK) \neq 2$ (as assumed throughout this section). Then the trivial representation of $H$ over $\bL$, has linear type \textbf{(A, A, 1, 1)}, and therefore type \textbf{(A, A, A, 1, 1, 2)}.

\subsection{Realisation of Type (2) / III}

Let $\bL / \bK$ be the extension $\bQ(i) / \bQ$, and $G = D_8 = \langle a,b \mid a^4 = b^2 = e, bab = a^{-1} \rangle$. Let $H$ be the subgroup $\{e, b, a^2, a^2 b\} \cong K_4$, and let $V$ be the $1$-dimensional representation of $H$ over $\bQ(i)$ with character $\chi$, where $\chi(b) = -1$ and $\chi(a^2) = -1$. Then the linear type of $V$ is \textbf{(A, A, 2, 2)}, and so the type of $V$ is \textbf{(A, A, B, 2, 2, 3(b))}.

\subsection{Realisation of Type (3)}

Let $\bL / \bK$ be the extension $\bQ(i) / \bQ$, and let $H = C_3$ inside $G = D_6$. Let $V$ be the non-trivial irreducible representation of $H$ over $\bQ(i)$, which is $2$-dimensional. The linear type of $V$ is \textbf{(A, A, 3(a), 3(a))}, and so the type of $V$ is \textbf{(A, A, A, 3(a), 3(a), 3(b))}.

\subsection{Realisation of Type (4) / II.V}

Let $\bL = \bQ(\zeta)$, where $\zeta$ is a primitive $8$th root of $1$. This is a degree $4$ extension of $\bQ$, and we take $\bK$ to be the intermediate field $\bQ(i)$. As $\bL = \bK(\zeta)$, with $\zeta^2 = i$, the non-trivial automorphism $\alpha$ of $\bL/\bK$ sends $\zeta \mapsto - \zeta$. Let $H = C_4 = \langle x \rangle$ and $G = C_8 = \langle y \rangle$, with $y^2 = x$. Let $W$ be the $1$-dimensional $\bK[H]$-module where $x$ acts by multiplication on $\bK$ by $i$, and let $V$ be the base change to $\bL$.

We claim that $V$ is extendable to a $\bL \rtimes G$-module. Define an $\alpha$-linear action of $y$ by
\[
    \rho_y(v) = (\zeta i) \cdot \alpha(v),
\]
for $v \in V$. Then we can compute that
\begin{align*}
\rho_y(\rho_y(v)) &= \rho_y((\zeta i)\alpha(v)), \\
&= (\zeta i) \alpha(\zeta i) \cdot \alpha^2(v), \\
&= \zeta (-\zeta) i^2 \cdot v, \\
&= \zeta^2 \cdot v, \\
&= i \cdot v,
\end{align*}
and therefore, as this is exactly the action of $x = y^2$, $V$ extends via this action of $y$ to an $\bL \rtimes G$-module.

To see that this example is of type \textbf{(B, A, A, 1, 3(b), 3(a))}, from the fact that $V$ is induced from a $\bK[H]$-module and extends to a $\bL \rtimes G$-module, it must be of type \textbf{(?, A, A, ?, ?, ?)}. The only other type with these specified entries is \textbf{(A, A, A, 1, 1, 2)}. In our case, we also have that $W$ is not the restriction of any irreducible $\bK[G]$-module, and therefore the extension $\bK[H] \subset \bK[G]$ must be in type 2 or 3(b), and so this example must have type \textbf{(B, A, A, 1, 3(b), 3(a))}.

\subsection{Realisation of Type (5) / II}

Let $\bL / \bK$ be the extension $\bQ(i) / \bQ$, $G = C_4$, and $H = C_2$. Then the linear type of the non-trivial $1$-dimensional $L$-representation $V$ is \textbf{(B, A, 1, 3(b))}. Therefore, to show that the type of $V$ is \textbf{(B, A, C, 1, 3(b), 3(b))}, it is sufficient to show that $V$ does not extend to a semilinear representation of $G$ over $L$.

Suppose there was such an extension, call it $\rho$. Writing $C_4 = \langle y \rangle$, $C_2 = \langle x \rangle$ with $y^2 = x$, let $a \in \bQ(i)$ be such that $y$ acts by matrix $a$. Then, from the relation $y^2 = x$, we have that
\[
a \overline{a} = -1
\]
which is impossible, as $a \overline{a} \in \bR_{\geq 0}$.

Alternatively, if such an extension of $V$ to a semilinear representation existed, it would also exist for $\bC / \bR$, which is impossible as the Frobenius-Schur indicator of $V$ is $-1$ (cf.\ Chapter \ref{chapter:tenfoldway}).

\subsection{Realisation of Type (6)}

Let $\bL / \bK$ be the extension $\bQ(\zeta_{21}) / \bQ(\zeta_7)$, let $G$ be the quaternion group $Q_8 = \langle a,x \mid a^4, x^2 = a^2, x^{-1}ax = a^{-1} \rangle$, and let $H = \langle a \rangle$. Let $V$ be the unique $2$-dimensional irreducible representation of $H$ over $\bQ(\zeta_{21})$, which has linear type \textbf{(C, A, 3(a), 3(b))}. To see that the type of $V$ is \textbf{(C, A, A, 3(a), 3(b), 3(a))} it is sufficient to show the following. 

\begin{lemma}
    There is an extension of $V$ to a semilinear representation of $G$ over $\bQ(\zeta_{21})$.
\end{lemma}

\begin{proof}
If $\mu \in \bQ(\zeta_7)$ and $a \in \bQ(\zeta_{21})$ with $N_{\bQ(\zeta_{21}) / \bQ(\zeta_{7})}(a)(1+\mu^2) = -1$, then the matrices
    \[
A = \begin{pmatrix}
    0 & -1 \\ 1 & 0
\end{pmatrix}, \qquad X = a \begin{pmatrix}
    1 & \mu \\ \mu & -1
\end{pmatrix}
\]
define such an extension of $V$, as $A^4 = I_2$, $X \cdot {}^{\sigma \!}X = - I_2 = A^2$, and $A \cdot X = X \cdot {}^{\sigma \!}A^{-1}$, where $\sigma$ denotes the non-trivial element of $\Gal(\bQ(\zeta_{21}) / \bQ(\zeta_{7}))$. Such a pair of elements $(\mu, a)$ can indeed be found: for example, taking $\mu = 2 \sqrt{-7} = 2(\zeta_7 + \zeta_7^2 - \zeta_7^3 + \zeta_7^4 - \zeta_7^5 - \zeta_7^6)$, and $a = \frac{1}{9}(\zeta_{21}^7 + 2)$.
\end{proof}

\subsection{Realisation of Type (7)}

Let $\bL / \bK$ be the extension $\bQ(\omega) / \bQ$, where $\omega$ is a primitive third root of $1$, and let $G$ be the quaternion group $Q_8 = \langle a,x \mid a^4, x^2 = a^2, x^{-1}ax = a^{-1} \rangle$. Taking $H = \langle a \rangle$, let $V$ be the unique $2$-dimensional irreducible representation of $H$ over $\bQ(\omega)$, which has linear type \textbf{(C, A, 3(a), 3(b))}. The following shows that the type of $V$ is \textbf{(C, A, C, 3(a), 3(b), 3(b))}. 

\begin{lemma}
    There is no extension of $V$ to a semilinear representation of $G$ over $\bQ(\omega)$.
\end{lemma}

\begin{proof}
A matrix form of $V$ over $\bQ(\omega)$ is given by 
\[
\rho(a) = \begin{pmatrix}
    0 & -1 \\ 1 & 0
\end{pmatrix},
\]
and an extension of $\rho$ to a semilinear representation of $G$ over $\bQ(\omega)$ is equivalent to the choice of an $X \in \GL_2(\bQ(\omega))$ with $A X = X \overline{A}^{-1}$ and $X \cdot \overline{X} = A^2 = -I_2$. The first condition implies that
\[
X = \begin{pmatrix}
    a & b \\ b & -a
\end{pmatrix}
\]
for some $a , b \in \bQ(\omega)$. The second condition is equivalent to $\overline{X} = - X^{-1}$, which mean that $a = \lambda \overline{a}$, $b = \lambda \overline{b}$, where $\lambda = \det(X) = - (a^2 + b^2)$. If $a \neq 0$, then $b/a = \lambda \overline{b} / \lambda \overline{a} = \overline{b/a}$, and similarly if $b \neq 0$ then $a/b \in \bQ$, so in either case there are $r,s \in \bR$ and $z \in S^1$ with $a = rz$ and $b = sz$. We have that $\lambda = -(a^2 + b^2) = - z^2 (r^2 + s^2)$, and from the equation $a = \lambda \overline{a}$ and the fact that $\overline{z} = z^{-1}$,
\[
r z = \lambda \overline{a} = - z^2 (r^2 + s^2) r \overline{z} = - r z (r^2 + s^2),
\]
so $r = 0$, as $r^2 + s^2 > 0$ because $(r,s) \in \bR^2 \setminus \mathbf{0}$. Similarly, this also forces $s = 0$, a contradiction.
\end{proof}

\subsection{Realisation of Type (8)}

Let $\bL / \bK$ be the extension $\bQ(i) / \bQ$, and $G = D_{16} = \langle a,b \mid a^8 = b^2 = e, bab = a^{-1} \rangle$. Let $H$ be the subgroup $\langle a^2, b \rangle \cong D_4$, and let $V$ be the unique $2$-dimensional representation of $H$ over $\bQ(i)$. Then the linear type of $V$ is \textbf{(A, A, 3(b), 3(b))}. To prove that the type of $V$ is \textbf{(A, A, A, 3(b), 3(b), 3(a))} it is sufficient to show that $V$ extends to a semilinear representation of $G$ over $\bQ(i)$. To this end, note that $V$ has a matrix representation $\rho$ with
\[
\rho(a^2) = \begin{pmatrix}
    i & 0 \\ 0 & -i
\end{pmatrix}, \qquad \rho(b) = \begin{pmatrix}
    0 & 1 \\ 1 & 0
\end{pmatrix}
\]
and so we may define
\[
A \coloneqq \begin{pmatrix}
    0 & -1 \\ i & 0
\end{pmatrix}
\]
which is chosen to satisfy the relations $A \overline{A} = \rho(a^2)$ and $A \overline{X} \overline{A} = X$. These are exactly the relations to ensure that $\rho(a) \coloneqq A$ defines an extension of $V$ to a semilinear representation of $G$ over $\bQ(i)$.

\subsection{Realisation of Type (9)} Let $G = C_{12}$, $H = C_6$, and $\bL / \bK$ be the extension $\bQ(\omega) / \bQ$ where $\omega$ is a primitive third root of $1$. Let $V$ be the unique order two character of $H$ over $\bQ(\omega)$. The linear type of $V$ is \textbf{(A, A, 3(b), 3(b))}, and to prove that the type of $V$ is \textbf{(A, A, C, 3(b), 3(b), 3(b))} it is sufficient to show that $V$ does not extend to a semilinear representation of $G$ over $\bQ(\omega)$.

Similarly to the realisation of type (5), such an extension is equivalent to the existence of an $a \in \bQ(\omega)$ with $a \overline{a} = -1$, which is impossible as the norm of an element of $\bQ(\omega)$ lies in $\bQ_{\geq 0}$, $\bQ(\omega) / \bQ$ being imaginary quadratic.

\subsection{Realisation of Type (10) / IV}

Let $H = C_3$, $G = C_6$, and $\bL/\bK$ be the extension $\bQ(\omega) / \bQ$, where $\omega$ is a primitive third root of $1$. Let $V$ be the $1$-dimensional representation of $H$ over $\bQ(\omega)$ where the generator acts by multiplication by $\omega$. Then the linear type of $V$ is \textbf{(B, B, 1, 1)}, and so the type of $V$ is \textbf{(B, B, B, 1, 1, 2)}.

\subsection{Realisation of Type (11) / VII}

Let $\bL / \bK$ be the extension $\bQ(i) / \bQ$, and let $G$ denote the modular group of order 16, identified as SmallGroup(16,6), which has presentation given by
\[
    G = \langle a, x \mid a^8 = x^2 = e, xax = a^5 \rangle.
\]  
This has a unique subgroup $H$ isomorphic to $C_4 \times C_2$, explicitly defined as $H = \langle a^2, x \rangle$. Let $V$ be the $1$-dimensional representation of $H$ over $\bQ(i)$ with character $\chi$ such that $\chi(a^2) = i$ and $\chi(x) = 1$. Then it is direct to check that the three characters $\chi, {}^a \chi, \overline{\chi}$ are distinct, and therefore $\Stab_{K_4}(V) = 1$, so from Table \ref{tab:24types5} the type of $V$ must be \textbf{(B, B, B, 2, 2, 1)}.

\subsection{Realisation of Type (12) / V}

Let $H = C_3$, $G = D_6$, and $\bL/\bK$ be the extension $\bQ(\omega) / \bQ$, where $\omega$ is a primitive third root of $1$. Let $V$ be the $1$-dimensional representation of $H$ over $\bQ(\omega)$ where the generator acts by multiplication by $\omega$. Then the linear type of $V$ is \textbf{(A, B, 2, 3(a))}, and so the type of $V$ is \textbf{(A, B, A, 2, 3(a), 1)}.

\subsection{Realisation of Type (13)}

Let $\bL / \bK$ be the extension $\bQ(i) / \bQ$, and let $G$ by the modular group of order 16 with presentation
\[
G = \langle a, x \mid a^8, x^2, xax^{-1} = a^5 \rangle.
\]
We take $H$ to be the cyclic subgroup $H = \langle a \rangle$, and $V$ to be either of the irreducible $2$-dimensional representations of $H$ over $\bQ(i)$. The linear type of $V$ is \textbf{(B, B, 3(a), 3(a))}, and so the type of $V$ is \textbf{(B, B, B, 3(a), 3(a), 3(b))}.

\subsection{Realisation of Type (14)}

Let $\bL / \bK$ be the extension $\bQ(i) / \bQ$, $G = C_8$ and $H = C_4$. Take $V$ to be either of the faithful characters of $H$ over $\bQ(i)$. The linear type of $V$ is \textbf{(B, B, 3(b), 3(b))}, and so the type of $V$ is \textbf{(B, B, B, 3(b), 3(b), 3(a))}.

\subsection{Realisation of Type (15) / VI}

Let $\bL / \bK$ be the extension $\bQ(i) / \bQ$, $G = Q_8$, and $H = \langle j \rangle \leq Q_8$. Take $V$ to be $1$-dimensional representation of $H$ with $\chi(j) = i$. The linear type of $V$ is \textbf{(C, B, 2, 3(b))}, and so the type of $V$ is \textbf{(C, B, C, 2, 3(b), 1)}.

\subsection{Realisation of Type (16) / VIII}

Let $\bL / \bK$ be the extension $\bQ(i) / \bQ$, $H = Q_8 \leq G = Q_8 \times C_2$. Take $V$ to be the unique $2$-dimensional representation of $H$. The linear type of $V$ is \textbf{(C, C, 1, 1)}, and so the type of $V$ is \textbf{(C, C, C, 1, 1, 2)}.

\subsection{Realisation of Type (17) / X}

Let $\bL / \bK$ be $\bQ(i) / \bQ$, and let $G$ be SmallGroup(32,8) in the database of small groups. The group $G$ has presentation
\[
G \;=\;
\left\langle 
  a, b, c \;\middle|\;
  a^{8} = b^{4} = c^{2} = 1,\;
  a^{4} = b^{2},\;
  bc = cb,\;
  a b a^{-1} = b^{3} c,\;
  a c a^{-1} = b^{2} c
\right\rangle
\]
and the field $\bQ(i)$ is a splitting field for $G$. We take $H$ to be the subgroup
\[
H = \langle a^2 , b, c \rangle
\]
which is isomorphic to $Q_8 \times C_2$, where, for $i$ and $j$ the standard generators of $Q_8$, $a^2 \leftrightarrow i$ and $b \leftrightarrow j$.

Let $V$ be the inflation of the unique $2$-dimensional representation of $Q_8$ along the projection $H \twoheadrightarrow Q_8$, tensored with the sign character of the projection $H \twoheadrightarrow C_2$. In particular, we have that $\chi(c) = -2$, $\chi(b^2) = -2$, and $\chi(b^2 c) = 2$. This character is $\bQ$-valued, but not induced from a representation over $\bQ$, and therefore the classical Schur index $m_c(V) = 2$. We can also check that ${}^{\omega} V \not\cong \overline{V}$, so $\Stab_{K_4}(V) = \langle a \rangle$, as $\chi(c) = -2$, but 
\[
{}^{a}\chi(c) = \chi(aca^{-1}) = \chi(b^2 c) = 2.
\] 
Consulting Table \ref{tab:24types5}, these two properties force the type of $V$ to be \textbf{(C, C, B, 2, 2, 3(a))}.

\subsection{Realisation of Type (18) / IX}

Let $\bL / \bK$ be the extension $\bQ(i) / \bQ$, and let $G$ be SmallGroup(16,13), also know as the Pauli group, with presentation
\[
G = \langle a, x, y \mid a^4, y^4, x^2, a^2 = y^2, xax^{-1} = a^{-1}, [a,y], [x,y] \rangle.
\]
There is a unique copy of $Q_8$ as the subgroup $H \coloneqq \langle a, xy \rangle$, where $i \mapsto a, j \mapsto xy$. Let $(V, \rho)$ be the unique $2$-dimensional irreducible representation of $H$, with matrix form
\[
\rho(a) = \begin{pmatrix}
    i & 0 \\ 0 & -i
\end{pmatrix}, \qquad \rho(xy) = \begin{pmatrix}
    0 & -1 \\ 1 & 0
\end{pmatrix}.
\]
This has linear type \textbf{(B, C, 1, 3(a))}. From Table \ref{tab:24types}, to show this has type \textbf{(B, C, A, 1, 3(a), 3(a))} it is sufficient to show that $\rho$ extends to a semilinear representation of $G$ over $\bQ(i)$. Set
\[
A \coloneqq \rho(a), \qquad X \coloneqq \begin{pmatrix}
    i & 0 \\ 0 & i
\end{pmatrix}, \qquad Y \coloneqq \begin{pmatrix}
    0 & -i \\ i & 0
\end{pmatrix}
\]
Then letting the generators $a, x, y$ of $G$ act by the matrices $A,X,Y$ respectively, one can check that this defines a semilinear representation which extends $\rho$, by explicitly verifying the relations
\[
\rho(a) = A, \ \rho(xy) = X \overline{Y} , \ Y \overline{Y} = A^2, \ X \overline{X} = I_2, \ X \overline{A} = A^{-1} X, \ A Y = \overline{Y} A, \ X \overline{Y} = Y \overline{X}.
\]

\subsection{Realisation of Type (19) / IX.V}\label{sect:realiseIXV} Let $\bL / \bK$ be the extension $\bQ(i) / \bQ$, and let $G$ be $\text{SmallGroup}(24,5)$, with presentation
\[
G = \langle w, x, y, z \;|\;
w^2, \;
x^2 = y, \;
y^2, \;
z^3, \;
[w,x], \;
[w,y], \;
[x,y], \;
[y,z], \;
[x,z], \;
w z w^{-1} = z^{-1}
\rangle.
\]
Let $H = \langle yz, wx \rangle$, which is isomorphic to
\[
\text{Dic}_{12} \coloneqq \langle a,b \mid a^6, b^2 = a^3,  b^{-1} a b = a^{-1} \rangle
\]
where $a$ corresponds to $yz$ and $b$ corresponds to $wx$. Let $V$ be the $2$-dimensional representation of $H$ over $\bQ(i)$ defined by 
\[
\rho(yz) \coloneqq \begin{pmatrix}
    0 & -1 \\ 1 & 1
\end{pmatrix}
\qquad
\rho(wx) \coloneqq \begin{pmatrix}
    0 & i \\ i & 0
\end{pmatrix}.
\]

We could compute the linear type of $V$ using the algorithm (cf.\ Section \ref{sect:lineartype}), but we instead give a more involved example than the previous types to show how this can be computed by hand.

\begin{lemma}
    The linear type of $V$ is \emph{\textbf{(B, C, 1, 3(a))}}.
\end{lemma}

\begin{proof}
The complex character table of $H$ is
\[
\begin{array}{c|c|rrrrrr}
 & m(\psi_i) & (1) & (a^3) & (a^2) & (b) & (ab) & (a) \\
\hline
\text{Size}  &   & 1 & 1 & 2 & 3 & 3 & 2 \\
\text{Order} &   & 1 & 2 & 3 & 4 & 4 & 6 \\
\hline
\psi_{1} & 1 & 1 & 1 & 1 & 1 & 1 & 1 \\
\psi_{2} & 1 & 1 & 1 & 1 & -1 & -1 & 1 \\
\psi_{3} & 1 & 1 & -1 & 1 & i & -i & -1 \\
\psi_{4} & 1 & 1 & -1 & 1 & -i & i & -1 \\
\psi_{5} & 2 & 2 & -2 & -1 & 0 & 0 & 1 \\
\psi_{6} & 1 & 2 & 2 & -1 & 0 & 0 & -1 \\
\end{array}
\]
where $m(\psi)$ is the Schur index of $\psi$ over $\bQ$. This is the same as the character table of $H$ over $\bQ(i)$: for $\psi_i$ with $i \neq 5$ this follows from the fact that $\psi_i$ is $\bQ(i)$-valued, and for $\psi_5$ this is because $\psi_5$ is the character of the representation $\rho$, and so admits a model over $\bQ(i)$. We know that the type of $\psi_5$ has C-type (the type of the extension $\bQ[H] \leq \bQ(i)[H]$) $\mathbf{C}$, and therefore that the type of $V$ is one of the nine types in the range (16) to (24) of Table \ref{tab:24types}. Furthermore, because $\psi_5$ is the unique representation with Schur index $2$, we see that conjugation by the odd element $w$ fixes $\psi_5$, and therefore by Table \ref{tab:indres} (which applies by Section \ref{sect:indextwosubgroup}) $V$ is the restriction of exactly two irreducible characters $\chi$, $\chi \otimes \pi$ of $G$ over $\bQ(i)$ with $\chi \neq \chi \otimes \pi$, where $\pi$ is the sign character of $G$ over $H$. Setting $Z \in \Irr_{\bQ(i)}(G)$ to be the representation corresponding to $\chi$, this immediately tells us that $b \not\in \Stab_{K_4}(Z)$, and therefore that the type of $V$ is one of (16), (18) or (19).

To see that the type of $V$ cannot be (16), consider the complex character table of $G$
\[
\begin{array}{c|rrrrrrrrrrrr}
\text{Class} & 1 & 2 & 3 & 4 & 5 & 6 & 7 & 8 & 9 & 10 & 11 & 12 \\
\hline
\text{Size}  & 1 & 1 & 3 & 3 & 2 & 1 & 1 & 3 & 3 & 2 & 2 & 2 \\
\text{Order} & 1 & 2 & 2 & 2 & 3 & 4 & 4 & 4 & 4 & 6 & 12 & 12 \\
\hline
\chi_1  & 1 & 1 & 1 & 1 & 1 & 1 & 1 & 1 & 1 & 1 & 1 & 1 \\
\chi_2  & 1 & 1 & -1 & -1 & 1 & 1 & 1 & -1 & -1 & 1 & 1 & 1 \\
\chi_3  & 1 & 1 & 1 & 1 & 1 & -1 & -1 & -1 & -1 & 1 & -1 & -1 \\
\chi_4  & 1 & 1 & -1 & -1 & 1 & -1 & -1 & 1 & 1 & 1 & -1 & -1 \\
\chi_5  & 1 & -1 & 1 & -1 & 1 & i & -i & i & -i & -1 & -i & i \\
\chi_6  & 1 & -1 & -1 & 1 & 1 & i & -i & -i & i & -1 & -i & i \\
\chi_7  & 1 & -1 & 1 & -1 & 1 & -i & i & -i & i & -1 & i & -i \\
\chi_8  & 1 & -1 & -1 & 1 & 1 & -i & i & i & -i & -1 & i & -i \\
\chi_9  & 2 & 2 & 0 & 0 & -1 & 2 & 2 & 0 & 0 & -1 & -1 & -1 \\
\chi_{10} & 2 & 2 & 0 & 0 & -1 & -2 & -2 & 0 & 0 & -1 & 1 & 1 \\
\chi_{11} & 2 & -2 & 0 & 0 & -1 & -2i & 2i & 0 & 0 & 1 & -i & i \\
\chi_{12} & 2 & -2 & 0 & 0 & -1 & 2i & -2i & 0 & 0 & 1 & i & -i \\
\end{array}
\]
The element $y$ is central in $G$, and order two, and therefore the 2nd column is the conjugacy class of $y$. We therefore see that the type characters $\chi$ and $\chi \otimes \pi$ which restrict to $\psi_5$ are $\chi_{11}$ and $\chi_{12}$, and so $\Stab_{K_4}(Z) = \langle c \rangle$ because $\overline{\chi_{11}} = \chi_{12}$. This tells us that the type of $V$ is one of (18) or (19), both of which have linear type \textbf{(B, C, 1, 3(a))}.
\end{proof}

Finally, to show that the type of $V$ is indeed \textbf{(B, C, C, 1, 3(a), 3(a))}, it remains to see that $V$ does not extend to a semilinear representation of $G$ over $\bQ(i)$.

\begin{lemma}\label{lem:noextension}
    There is no extension of $V$ to a semilinear representation of $G$ over $\bQ(i)$.
\end{lemma}
\begin{proof}
Suppose there exists such an extension. We have that $(yz)^3 = y$ and $(yz)^4 = z$, so $\rho(y) = -I_2$ and $\rho(z) = -\rho(yz)$. Furthermore, both generators $w$ and $x$ are odd (in $G \setminus H$), and therefore such an extension is equivalent to the data of matrices $W,X$ over $\bQ(i)$ which satisfy the appropriate relations coming from the relations in $G$ and $\rho(wx) = W \overline{X}$. In particular, the relations $w^2 = 1$, $x^2 = y$ and $xz = zx$ mean that $W$ and $X$ satisfy the relations
\begin{align*}
W \overline{W} &= I_2, \\ 
X \overline{X} &= \rho(y), \\
X \overline{\rho(z)} &= \rho(z) X.
\end{align*}
Writing
\[
X = \begin{pmatrix} a & b \\ c & d \end{pmatrix},
\]
then as $\rho(z)$ is $\bQ$-valued
\[
X \rho(z) = \rho(z) X,
\]
and therefore
\[
\begin{pmatrix} a & b \\ c & d \end{pmatrix} \begin{pmatrix} 0 & 1 \\ -1 & -1 \end{pmatrix} 
=
\begin{pmatrix} 0 & 1 \\ -1 & -1 \end{pmatrix}\begin{pmatrix} a & b \\ c & d \end{pmatrix}
\]
or equivalently
\[
\begin{pmatrix} -b & a-b \\ -d & c-d \end{pmatrix} = \begin{pmatrix} c & d \\ -a-c & -b-d \end{pmatrix}.
\]
In particular, $c = -b$ and $d = a-b$, so
\[
X = \begin{pmatrix} a & b \\ -b & a-b \end{pmatrix}.
\]
Because $b = wx$ and $w$ commutes with $x$, $x$ commutes with $b$ and therefore
\[
X \overline{\rho(b)} = \rho(b) X
\]
so
\[
\begin{pmatrix} a & b \\ -b & a-b \end{pmatrix} \begin{pmatrix} 0 & -i \\ -i & 0 \end{pmatrix}
=
\begin{pmatrix} 0 & i \\ i & 0 \end{pmatrix} \begin{pmatrix} a & b \\ -b & a-b \end{pmatrix}
\]
or equivalently
\[
\begin{pmatrix} -b & a-b \\ a & b \end{pmatrix} = \begin{pmatrix} -b & -a \\ b-a & b \end{pmatrix}.
\]
In particular $b -a =a$ so $b = 2a$ and
\[
X = \begin{pmatrix} a & 2a \\ -2a & -a \end{pmatrix}.
\]
Finally, we have that
\[
X \overline{X} = \rho(y) = -I_2,
\]
so $X^{-1} = - \overline{X}$, hence
\[
- \det(X) \begin{pmatrix} \overline{a} & \overline{2a} \\ -\overline{2a} & -\overline{a} \end{pmatrix} = \begin{pmatrix} -a & -2a \\ 2a & a \end{pmatrix}.
\]
Concentrating on any fixed entry,
\[
a = \overline{a} \det(X) = 3 a^2 \overline{a},
\]
hence
\[
\frac{1}{3} = a \overline{a}.
\]
However, this is impossible: the $3$-adic valuation of the norm of an element of $\bQ(i)$ is always even, because $3$ is inert in $\bQ(i)$.
\end{proof}

\begin{remark}
If $\bL$ was bigger then this representation over $\bK$ would extend. For example, if $\bL / \bK$ was $\bC / \bR$ we could take $a = 1 / \sqrt{3}$ and defining $X$ as above would give the required extension.
\end{remark}

\begin{remark}\label{rem:differentdivalg}
    Let us identify the division algebras $D_1 \coloneqq \End_{\bQ[H]}(X)$ and $D_2 \coloneqq \End_{\bQ(i) \rtimes G}(W)$ in $\Br(\bQ)$. Using the \texttt{SchurIndices} function in magma, one can check that the non-split places of $D_1$ are $3$ and $\infty$. For $D_2$, even though we do not know this algebra explicitly (we haven't even constructed explicitly the semilinear representation $W$), we can still identify $D_2$ by considering the splitting behaviour at places of $\bQ$.
    
    For a place $v$ of $\bQ$, $D_{2,v}$ is split whenever $\bQ(i)_v \coloneqq \bQ_v \otimes_{\bQ} \bQ(i)$ is not a field \cite[Cor.\ 11.4]{TAYSchur}. Suppose instead that $\bQ(i)_v$ is a field, which happens exactly when $v = \infty$ or $v = p$ is a finite prime $p \geq 2$ with $p = 2$ or $p \equiv 3 \!\! \mod 4$. Note that in this case $\bQ(i)_v = \bQ_v(i)$, and
    \[
    V_v \coloneqq \bQ(i)_v \otimes_{\bQ(i)} V \in \Rep_{\bQ(i)_v}(H)
    \]
    is irreducible, as $V$ is absolutely irreducible over $\bQ(i)$, and therefore $V_v$ has a Schur index $m_{\bQ_v}^{\bQ_v(i)}(V_v) \in \{1,2\}$. This Schur index is related to the splitting of $D_{2,v}$ by \cite[Cor.\ 11.4]{TAYSchur}, which tells us that $D_{2,v}$ is split if and only if $m_{\bQ_v}^{\bQ_v(i)}(V_v) = 1$. From the proof of Lemma \ref{lem:noextension}, one sees that $m_{\bQ_v}^{\bQ_v(i)}(V_v) = 2$ (equiv.\ $V_v$ does not extend to a semilinear representation of $G$ over $\bQ(i)_v$) if and only if $1/3$ cannot be written as the sum of two squares in $\bQ_v$. Over $\bR$, $1/3$ is the sum of squares, so $D_{2,v}$ is split over $\bR$ (and so in particular $D_1 \not\cong D_2$). 

    For the other places, if $p \geq 2$ is a finite prime then $1/3$ is the sum of two squares in $\bQ_p$ if and only if $3p^{2e}$ is the sum of the squares of two elements of $\bZ_p$ for some $e \geq 1$, which can be seen by clearing denominators. For primes $p \equiv 3 \!\! \mod 4$, a non-zero $p$-adic integer $x \in \bZ_p$ is a sum of two squares of $\bZ_p$ if and only if $\nu_p(x)$ is even, and a non-zero $2$-adic integer $x = 2^r u \in \bZ_2$ with $u \in \bZ_2^{\times}$ is a sum of two squares in $\bZ_2$ if and only if $u \equiv 1 \!\! \mod 4$ \cite{KCON}. In particular, taking $x \coloneqq 3p^{2e}$, we have that $3p^{2e}$ is not a sum of two elements of $\bZ_p$ if and only if $p = 2$ or $p = 3$.

    In summary, we can conclude that places over $\bQ$ where $D_2$ splits are the finite places $2$ and $3$. 
\end{remark}

\subsection{Realisation of Type (20)}

Let $\bL / \bK$ be the extension $\bQ(i) / \bQ$ and $G = \text{SmallGroup}(32,44)$, with presentation
\[
G = \langle x, y, z \mid x^8,\; y^2,\; z^2,\; yxy^{-1} = x^3,\; zxz^{-1} = x^5,\; zyz^{-1} = x^4 y\rangle.
\]
We take $H \coloneqq \langle x, yz \rangle$, which is isomorphic to $Q_{16}$, and let $V$ be the unique $4$-dimensional irreducible representation of $H$ over $\bQ(i)$, explicitly given by
\[
\rho(x) = \begin{pmatrix}
    0 & i & 0 & 0 \\
1 & 0 & 0 & 0 \\
0 & 0 & 0 & 1 \\
0 & 0 & -i & 0
\end{pmatrix}, \qquad \rho(yz) = \begin{pmatrix}
0 & 0 & -1 & 0 \\
0 & 0 & 0 & -1 \\
1 & 0 & 0 & 0 \\
0 & 1 & 0 & 0
\end{pmatrix}.
\]
The linear type of $V$ is \textbf{(C, C, 3(a), 3(a))}, and therefore to show that the type of $V$ is \textbf{(C, C, A, 3(a), 3(a), 3(a))} it is sufficient to show that $V$ extends to a semilinear representation of $G$ over $\bQ(i)$. Set
\[
Y = \begin{pmatrix}
-1 & 0 & 0 & 0 \\
0 & -i & 0 & 0 \\
0 & 0 & 1 & 0 \\
0 & 0 & 0 & i
\end{pmatrix}, \qquad Z = \begin{pmatrix}
0 & 0 & 1 & 0 \\
0 & 0 & 0 & i \\
1 & 0 & 0 & 0 \\
0 & i & 0 & 0
\end{pmatrix}.
\]
Then letting the generators $y$ and $z$ of $G$ act by the matrices $Y$ and $Z$ respectively, one can check that this defines a semilinear representation which extends $\rho$, by explicitly verifying the relations
\[
\rho(yz) = Y \overline{Z} , \quad Y \overline{Y} = I_2, \quad Z \overline{Z} = I_2, \quad Y \overline{\rho(x)} = \rho(x)^3 Y, \quad Z \overline{\rho(x)} = \rho(x)^5 Z, \quad Z \overline{Y} = \rho(x)^4 Y \overline{Z}. 
\]

\subsection{Analysis of Type (21)}

We do not know if type (21) can be realised. This has the same linear type as type (20), and we have checked by hand that all representations $V \in \Irr_{\bbL}(H)$ for $\bL / \bK$ either $\bQ(i) / \bQ$ or $\bQ(\omega) / \bQ$ and $|H| \leq 50$ with the correct linear type have type (20), not type (21). Therefore we leave the following as an open question.

\begin{prob}\label{prob:case21}
Is type (21) realisable, for some finite extension $\bL / \bK$ and finite group $G$?
\end{prob}

It is unclear if this is true even without assuming that $G$ is finite. Let us give a brief analysis of the behaviour of this type, to explain the constraints on the linear representations that it imposes, in the case when $V$ has minimal possible dimension $4$ over $\bL$.

In this case, $D_U \coloneqq \End(U)$ is a $4$-dimensional central division $\bK$-algebra. Let $D_X \coloneqq \End(X)$, and $D_W \coloneqq \End(W)$, which both have dimension $8$ over $\bK$. We have that $U|_H = X$, and $\bL_{\star \hspace{-0.1em}}[G] \otimes_{\bL[H]} U = W$, and so there are maps $D_U \rightarrow D_X$ and $D_U \rightarrow D_W$. We therefore have isomorphisms
\[
Z_X \otimes_{\bbK} D_U \xrightarrow{\sim} D_X, \qquad Z_W \otimes_{\bbK} D_U \xrightarrow{\sim} D_W,
\]
because these maps are injective (from central simple algebras) and the image has dimension $8$.

We have that $\End(V)$ is a degree four field extension of $\bK$. We also have that $X_L = V \oplus V$, and $W|_H = V \oplus V$, hence
\[
\bL \otimes_{\bbK} D_X \xrightarrow{\sim} M_2(\End(V)), \qquad \bL \otimes_{\bbK} D_W \xrightarrow{\sim} M_2(\End(V)),
\]
and taking the centre
\[
\bL \otimes_{\bbK} Z_X \xrightarrow{\sim} \End(V), \qquad \bL \otimes_{\bbK} Z_W \xrightarrow{\sim} \End(V).
\]

Consider the extension $\bL_X \coloneqq \bL \otimes_{\bbK} Z_X / Z_X$. Then
\[
\End(X_{Z_X}) \cong Z_X \otimes_{\bbK} D_X \cong Z_X \otimes_{\bbK} Z_X \otimes_{\bbK} D_U \cong (Z_X \times Z_X) \otimes_{\bbK} D_U = D_X \times D_X.
\]
Therefore $X_{Z_X} = X_1 \oplus X_2$, with $X_1 \not\cong X_2$ and $\End(X_1) \cong D_X \cong \End(X_2)$ as $Z_X$-algebras. Furthermore, as $X_1 \not\cong X_2$, then $V_1 \coloneqq X_{1, \bbL_X} \not\cong V_2 \coloneqq X_{2, \bbL_X}$. In particular,
\[
V_{\bbL_X} \cong V_1 \oplus V_2
\]
with $V_1 \not\cong V_2$. Now $Z$ is absolutely irreducible, and $V_1$ and $V_2$ both correspond under induction-restriction to $Z_{\bbL_X}$, and $Z \otimes \pi \cong Z$, hence $V_i \not\cong V_i^{\omega}$, and thus $V_1 \cong V_2^{\omega}$, as $V^{\omega} = V$. Furthermore $V_i \cong \overline{V_i}$, applying the non-trivial automorphism of $\bL_X / Z_X$ as each is defined over $Z_X$, and therefore $V_1 \not\cong \overline{V_1}^{\omega} \cong V_2$. Both $V_1$ and $V_2$ are absolutely irreducible, and therefore there is an irreducible semilinear representation $W_0 / \bL_X$ with $W_0|_H = V_1 \oplus V_2$ and $\End(W_0) = \bL_X$.

Now $W|_H = V \oplus V$, hence
\[
W_{\bbL_X}|_H = (W|_H)_{\bbL_X} = V_1 \oplus V_2 \oplus V_1 \oplus V_2.
\]
In particular, $W_{\bbL_X}|_H = (W_0 \oplus W_0)|_H$, hence $W_{\bbL_X} = W_0 \oplus W_0$, and
\[
Z_X \otimes_{\bbK} \End(W) \cong \End(W_{\bbL_X}) \cong M_2(\End(W_0)) \cong M_2(\bL_X).
\]
For example, taking the centre, we have that 
\[
Z_X \otimes_{\bbK} Z_W \xrightarrow{\sim} \bL_X.
\]
In particular, we see that $Z_X$ and $Z_W$ are linearly disjoint. They are two of the three index two subfields of $\bL_X = \End(V)$, the other being $\bL$. For example, this means that $Z_X \neq Z_W$.

\subsection{Realisation of Type (22)}

Let $\bL / \bK$ be the extension $\bQ(i) / \bQ$, and let $G = \text{SD}_{16}$, the semidihedral group of order 16, with presentation
\[
G = \langle a,x \mid a^8, x^2, xax^{-1} = a^3\rangle.
\]
We take $H \coloneqq \langle a^2, ax \rangle$, which is isomorphic to $Q_8$, and let $V$ be the unique $2$-dimensional irreducible representation of $H$ over $\bQ(i)$. The linear type of $V$ is \textbf{(A, C, 3(b), 3(a))}, and therefore to show that the type of $V$ is \textbf{(A, C, A, 3(b), 3(a), 3(a))} it is sufficient to show that $V$ extends to a semilinear representation of $G$ over $\bQ(i)$. The matrix form of $V$ is 
\[
\rho(a^2) = \begin{pmatrix}
    i & 0 \\ 0 & -i
\end{pmatrix}, \qquad \rho(ax) = \begin{pmatrix}
    0 & -1 \\ 1 & 0
\end{pmatrix}.
\]
Set
\[
A \coloneqq \begin{pmatrix}
    0 & -1 \\ i & 0
\end{pmatrix}, \qquad X \coloneqq \begin{pmatrix}
    i & 0 \\ 0 & 1
\end{pmatrix}.
\]
Then letting the generators $a$ and $x$ of $G$ act by the matrices $A$ and $X$ respectively, one can check that this defines a semilinear representation which extends $\rho$, by explicitly verifying the relations
\[
\rho(a^2) = A \overline{A} , \quad \rho(ax) = A \overline{X}, \quad X \overline{X} = I_2, \quad X \overline{A} X = A \overline{A} A.
\]

\subsection{Realisation of Type (23)}

Let $\bL / \bK$ be the extension $\bQ(\omega) / \bQ$, where $\omega$ is a primitive third root of $1$, and let $G = \text{SmallGroup}(16,13)$, with presentation
\[
G = \langle x,y,z \mid x^4, z^2, y^2 = z^2, [x,y], [x,z], zyz^{-1} = y^{-1}\rangle.
\]
This has a unique subgroup isomorphic to $Q_8$, given by $H = \langle y, xz\rangle$, where $y$ and $xz$ correspond to the generators $i$, $j$ of $Q_8$. Let $V$ be the $2$-dimensional representation irreducible representation of $Q_8$ over $\bQ(\omega)$. The linear type of $V$ is \textbf{(A, C, 3(b), 3(a))}, and therefore to show that the type of $V$ is \textbf{(A, C, C, 3(b), 3(a), 3(b))} it is sufficient to show the following.

\begin{lemma}\label{lem:noextension23}
    There is no extension of $V$ to a semilinear representation of $G$ over $\bQ(\omega)$.
\end{lemma}

\begin{proof}
This is similar to the proof of Lemma \ref{lem:noextension}. A matrix form of $V$ over $\bQ(\omega)$ is given by 
\[
\rho(y) = \begin{pmatrix}
    \omega & \omega^2 \\ \omega^2 & -\omega
\end{pmatrix}, \qquad
\rho(xy) = \begin{pmatrix}
    0 & -1 \\ 1 & 0
\end{pmatrix}
\]
and an extension of $\rho$ to a semilinear representation of $G$ over $\bQ(\omega)$ is equivalent to the data of a $X,Z \in \GL_2(\bQ(\omega))$ such that the matrix relations corresponding to the relations of $G$ are satisfied and $X\overline{Z} = \rho(xy)$. From the fact that $zy = y^{-1}z$, we have that $Z \overline{\rho(y)} = \rho(y^{-1}) Z$, from which one obtains that $Z$ has the form 
\[
Z = \begin{pmatrix}
    \omega b + \omega^2 c & b \\ c & -\omega^2 b - \omega c
\end{pmatrix}
\]
for some $b,c \in \bQ(\omega)$. From $z^2 = 1$, $\overline{Z} Z = I_2$, and from this one can compute that $b = -\det(Z)\overline{b}$ and $c = - \det(Z) \overline{c}$. From $x^2 = b^2$, $X \overline{X} = \rho(b^2) = - I_2$, and $X = \rho(xy)Z$ from $X\overline{Z} = \rho(xy)$, so $\rho(xy) Z \rho(xy) \overline{Z} = -I_2$. From this one can again compute that $b = \det(Z) \overline{c}$ and $c = \det(Z) \overline{b}$. Combining, we have that $c = -b$, and computing the determinant that $\det(Z) = -2 b^2$. Putting these back into $b = -\det(Z) \overline{b}$, $b = 2b^2\overline{b}$, so $1/2 = b \overline{b}$. However this is impossible: the $2$-adic valuation of a norm of an element of $\bQ(\omega)$ is always even, because the prime $2$ is inert in $\bQ(\omega)$.
\end{proof}

\subsection{Realisation of Type (24)}

Let $\bL / \bK$ be the extension $\bQ(i) / \bQ$, and let $G = Q_{16}$, the generalised quaternion group of order 16, with presentation
\[
G = \langle a,b,c \mid a^4 = b^2 = c^2 = abc \rangle.
\]
We take $H \coloneqq \langle a^2, b \rangle$, which is isomorphic to $Q_8$, and let $V$ be the unique $2$-dimensional irreducible representation of $H$ over $\bQ(i)$. The linear type of $V$ is \textbf{(C, C, 3(b), 3(b))}, and so the type of $V$ is \textbf{(C, C, C, 3(b), 3(b), 3(a))}.

\section{Interpretation of the 24 Types in Terms of Blocks}\label{sect:blockinterpretation}

We can interpret the 24 types of Theorem \ref{thm:24ways} in terms of extensions of matrix algebras coming from blocks of Diagram \eqref{eqn:diamondrings_complete}. From Theorem \ref{SBlock_3}, and $\bK$ has characteristic different from $2$, we know that the inclusions of Diagram \eqref{eqn:diamondrings} induce maps of Bohr completions, and the each resulting morphism of Diagram \eqref{eqn:diamondrings_complete} is injective.

Suppose that $e \in \widehat{\bK[H]}$ is a central primitive idempotent. There is no guarantee that the image of $e$ is central in any of the larger algebras of Diagram \eqref{eqn:diamondrings_complete}, so we consider
\[
\zeta_{\omega} \colon \widehat{\bK[H]} \rightarrow \widehat{\bK[H]},
\]
the $\bK$-algebra automorphism defined by completing the $\bK$-algebra automorphism
\[
\omega^{-1} - \omega \colon \bK[H] \rightarrow \bK[H]
\]
of $\bK[H]$, and set
\[
f \coloneqq \twopartdef{e}{e = \zeta_{\omega}(e)}{e + \zeta_{\omega}(e)}{e \neq \zeta_{\omega}(e)}.
\] 
Concretely, if $e$ corresponds to the simple module $X \in \Irr_{\bbK}(H)$ under the correspondence of Lemma \ref{lem:idemirrcorresp}, then $\zeta_{\omega}(e)$ corresponds to ${}^{\omega} X \in \Irr_{\bbK}(H)$.

For $B$ any of the $\bK$-algebras of Diagram \ref{eqn:diamondrings}, we write $\widehat{\iota}_B \colon \widehat{\bK[H]} \hookrightarrow \widehat{B}$ for the inclusion map.

\begin{defn}\label{def:ideminbigger}
    For any central primitive idempotent $e \in \widehat{\bK[H]}$, the \emph{semilinear block}, or \emph{S-block}, corresponding to $e$ is the diagram of $\bK$-algebra inclusions
    \[
     \begin{tikzcd}
	& \widehat{\iota}_{\bbK[G]}(f) \cdot \widehat{\bK[G]} \\
	\widehat{\iota}_{\bbL[G]}(f) \cdot \widehat{\bL[G]} & f \cdot \widehat{\bK[H]} & \widehat{\iota}_{\bL_{\star \hspace{-0.1em}}[G]}(f) \cdot \widehat{\bL_{\star \hspace{-0.1em}}[G]} \\
	& \widehat{\iota}_{\bbL[H]}(f)\cdot \widehat{\bL[H]}
	\arrow[hook, from=1-2, to=2-1]
	\arrow[hook, from=1-2, to=2-3]
	\arrow[hook, from=2-2, to=1-2]
	\arrow[hook, from=2-2, to=3-2]
	\arrow[hook, from=3-2, to=2-1]
	\arrow[hook, from=3-2, to=2-3]
\end{tikzcd}
\]
\end{defn}

We now interpret the 24 types of Theorem \ref{thm:24ways} in terms of the S-block. In particular, this shows that the image of $e$ in each larger algebra is a central idempotent, and that each algebra of the S-block is a matrix algebra.

\begin{prop}\label{prop:descofSblock}
    Suppose that $V \in \Irr_{\bbL}(H)$, and that $U$, $W$, $X$ and $Z$ correspond to $V$ as in Section \ref{sect:propagation}. Let $e \in \widehat{\bK[H]}$ be the central primitive idempotent corresponding to $X$, and define $f \in \widehat{\bK[H]}$ as above. Then:
    \begin{itemize}
        \item $f = \sum_{T \in C_2 \cdot X} e_T$, and for $n_X \coloneqq \dim_{\bbK}(X) / \dim_{\bbK}(\End_{\bbK[H]}(X))$,
        \[
            f \cdot \widehat{\bK[H]} \cong \prod_{T \in C_2 \cdot X} M_{n_X}(\End_{\bbK[H]}(X)).
        \]
        \item $\widehat{\iota}_{\bbK[G]}(f) = \sum_{T \in C_2 \cdot U} e_T$, and for $n_U \coloneqq \dim_{\bbK}(U) / \dim_{\bbK}(\End_{\bbK[G]}(U))$,
        \[
            \widehat{\iota}_{\bbK[G]}(f) \cdot \widehat{\bK[G]} \cong \prod_{T \in C_2 \cdot U} M_{n_U}(\End_{\bbK[G]}(U)).
        \]
        \item $\widehat{\iota}_{\bbL[H]}(f) = \sum_{T \in K_4 \cdot V} e_T$, and for $n_V \coloneqq \dim_{\bbL}(V) / \dim_{\bbL}(\End_{\bbL[H]}(V))$,
        \[
            \widehat{\iota}_{\bbL[H]}(f) \cdot \widehat{\bL[H]} \cong \prod_{T \in K_4 \cdot V} M_{n_V}(\End_{\bbL[H]}(V)).
        \]
        \item $\widehat{\iota}_{\bbL_{\star \hspace{-0.1em}}[G]}(f) = \sum_{T \in C_2 \cdot W} e_T$, and for $n_W \coloneqq \dim_{\bbK}(W) / \dim_{\bbK}(\End_{\bbL_{\star \hspace{-0.1em}}[G]}(W))$,
        \[
            \widehat{\iota}_{\bbL_{\star \hspace{-0.1em}}[G]}(f) \cdot \widehat{\bL_{\star \hspace{-0.1em}}[G]} \cong \prod_{T \in C_2 \cdot W} M_{n_W}(\End_{\bbL_{\star \hspace{-0.1em}}[G]}(W)).
        \]
        \item $\widehat{\iota}_{\bbL[G]}(f) = \sum_{T \in K_4 \cdot Z} e_T$, and for $n_Z \coloneqq \dim_{\bbL}(Z) / \dim_{\bbL}(\End_{\bbL[H]}(V))$,
        \[
            \widehat{\iota}_{\bbL[G]}(f) \cdot \widehat{\bL[G]} \cong \prod_{T \in K_4 \cdot Z} M_{n_Z}(\End_{\bbL[G]}(Z)).
        \]
    \end{itemize}
\end{prop}

\begin{proof}
We give the argument for $\bK[G]$, the remaining points being deduced in exactly the same manner. Note that $\widehat{\iota}_{\bbK[G]}(f)$ acts as the identity on simple objects of $\Rep_{\bbK}(G)$ isomorphic to $U$ or $U \otimes \pi$, because these restrict in $\Rep_{\bbK}(H)$ to a direct sum of $X$ and ${}^{\omega}X$, and as zero on all other simple representations of $\bK$ over $G$, as these contain no copies of $X$ or ${}^{\omega}X$ in their restriction to $\Rep_{\bbK}(H)$. This exactly means that $\widehat{\iota}_{\bbK[G]}(f)$ is the sum of $e_{U}$ and $e_{U \otimes \pi}$, the central primitive idempotents corresponding to $X$ and ${}^{\omega} X$ respectively under the correspondence of Lemma \ref{lem:idemirrcorresp}. The second statement follows immediately, from the definition of the idempotents and the description of the Bohr completion as the product of matrix rings over division rings of Lemma \ref{Bohr2}, noting that the endomorphism rings of both $U$ and $U \otimes \pi$ are the same. 
\end{proof}

\begin{eg}
Suppose that $V$ is the trivial representation of $H$. Then endomorphism rings of the corresponding irreducible representations are listed in Table \ref{tab:12types}, the orbits in Table \ref{tab:24types5}, and the dimensions in Table \ref{tab:24types5.5}, and thus by Proposition \ref{prop:descofSblock} the S-block is
\[
     \begin{tikzcd}
	& \bK \times \bK &  \\
	\bL \times \bL & \bK & M_2(\bK) \\
	& \bL &
	\arrow[hook, from=1-2, to=2-1]
	\arrow[hook, from=1-2, to=2-3]
	\arrow[hook, from=2-2, to=1-2]
	\arrow[hook, from=2-2, to=3-2]
	\arrow[hook, from=3-2, to=2-1]
	\arrow[hook, from=3-2, to=2-3]
\end{tikzcd}
\]
\end{eg}

\begin{eg}
For a more interesting example, suppose that $V$ is absolutely irreducible of type VI $ = (15)$, of dimension $n$. As above, consulting Tables \ref{tab:24types5}, \ref{tab:24types5.5} and \ref{tab:12types}, the S-block for $V$ is
\[
     \qquad \qquad \begin{tikzcd}
	& M_n(D) \\
	M_{2n}(\bL) & M_n(\bL) & M_n(D) \times M_n(D) \\
	& M_n(\bL) \times M_n(\bL)
	\arrow[hook, from=1-2, to=2-1]
	\arrow[hook, from=1-2, to=2-3]
	\arrow[hook, from=2-2, to=1-2]
	\arrow[hook, from=2-2, to=3-2]
	\arrow[hook, from=3-2, to=2-1]
	\arrow[hook, from=3-2, to=2-3]
\end{tikzcd}
\]
where $D$ is a $4$-dimensional central division algebra over $\bK$.
\end{eg}
 
\section{Associated $C_2$-Graded Algebras}\label{sect:BrauerWall}

In this section we associate to any semilinear representation $W \in \Irr_{\bbL}^{\hspace{0.1em} \star}(G)$ a $C_2$-graded division algebra $D(W)$, and a simple $C_2$-graded algebra $A(W)$, and explore how these put more restrictions on the structure of the S-block, or equivalently the type of $W$.

For basic notions concerning $C_2$-graded algebras, we direct the reader to Appendix \ref{appendix:gradedalgebras}, which contains the main definitions and certain results which we will make use of below.

We first consider $D(W)$.

\begin{defn}
If $W \in \Irr_{\bbL}^{\hspace{0.1em} \star}(G)$, we set
\[
D(W) \coloneqq \End_{\bbK[G]}(W),
\]
with $C_2$-grading
\begin{align*}
D(W)_0 &= \{f \in \End_{\bbK[G]}(W) \mid f(\lambda w) = \lambda f(w) \text{ for any $w \in W$, $\lambda \in \bL$} \}, \\
D(W)_1 &= \{f \in \End_{\bbK[G]}(W) \mid f(\lambda w) = \overline{\lambda} f(w) \text{ for any $w \in W$, $\lambda \in \bL$}\}.
\end{align*}
\end{defn}

To see that this is a well-defined $C_2$-grading, meaning that $D(W) = D(W)_0 \oplus D(W)_1$, take $\alpha \in \bL \setminus \bK$ with $\overline{\alpha} = - \alpha$, which exists from the eigenspace decomposition for $\overline{(-)} \colon \bL \rightarrow \bL$,
\[
\bL = \{\alpha \in \bL \mid \overline{\alpha} = \alpha\} \oplus \{\alpha \in \bL \mid \overline{\alpha} = -\alpha\},
\]
which uses that $\overline{(-)} \colon \bL \rightarrow \bL$ has order two and $\charfield(\bK) \neq 2$. We can define the automorphism
\[
t \coloneqq t_{\alpha} \colon \End_{\bbK[G]}(W) \rightarrow \End_{\bbK[G]}(W), \qquad t_{\alpha}(f)(w) = \frac{1}{\alpha} f(\alpha w),
\]
which has order two because $\alpha^2 \in \bK$, and is independent of the choice of $\alpha$.

Then, just as with $\bL$, because $\charfield(\bK) \neq 2$ we can decompose $D(W)$ into $\pm 1$ eigenspaces for $t_{\alpha}$, which are exactly $D(W)_0$ and $D(W)_1$. Explicitly, given $f \in D(W)$,
\[
f = f_0 + f_1, \qquad f_0 \coloneqq \frac{\alpha f(v) + f(\alpha v)}{2 \alpha} \in D(W)_0, \qquad f_1 \coloneqq \frac{\alpha f(v) - f(\alpha v)}{2 \alpha} \in D(W)_1.
\]

To see this more canonically, note that there is a $\bL_{\star \hspace{-0.1em}}[G]$-linear isomorphism
\[
\bL \otimes_{\bbK} W \xrightarrow{\sim} W \oplus \overline{W}, \qquad \lambda \otimes w \mapsto (\lambda w, \overline{\lambda}w)
\]
as in the proof of Proposition \ref{prop:resindKGLG} (which incidently doesn't require that $\charfield(\bK) \neq 2$), and thus we can compose the isomorphisms
\begin{align*}
\Hom_{\bbL_{\star \hspace{-0.1em}}[G]}(W,W) \oplus \Hom_{\bbL_{\star \hspace{-0.1em}}[G]}(\overline{W},W) \xrightarrow{\sim} \Hom_{\bbL_{\star \hspace{-0.1em}}[G]}(\bL \otimes_{\bbK} W , W) \xrightarrow{\sim} \Hom_{\bbK[G]}(W,W),
\end{align*}
where
\[
f_1 \oplus f_2 \mapsto [w \mapsto f_1(w) + f_2(w)],
\]
under which $\Hom_{\bbL_{\star \hspace{-0.1em}}[G]}(W,W)$ has image $D(W)_0$ and $\Hom_{\bbL_{\star \hspace{-0.1em}}[G]}(\overline{W},W)$ has image $D(W)_1$. This perspective allows us to immediately deduce the following basic properties of $D(W)$.

\begin{prop}\label{prop:basicpropDW}
    For $W \in \Irr_{\bbL}^{\hspace{0.1em} \star}(G)$, $D(W)$ satisfies the following properties:
    \begin{enumerate}
        \item $D(W)$ is a $C_2$-graded division algebra,
        \item $D(W)$ is ungraded \emph{(}meaning $D(W)_1 = 0$\emph{)} if and only if $W \not\cong \overline{W}$,
        \item $D(W)$ is simple as a $\bK$-algebra if and only if $U \cong U \otimes \pi$,
        \item The graded centre of $D(W)$ is contained in $D(W)_0$, so
        \[
            \widehat{Z}(D(W)) = D(W)_0 \cap Z(D(W)) \leq Z(D(W)_0).
        \] 
    \end{enumerate}
\end{prop}

\begin{proof}
    The first two properties follow immediately from the above discussion and Schur's lemma, noting that the inverse of any non-zero element of $x \in \Hom_{\bbL_{\star \hspace{-0.1em}}[G]}(\overline{W},W)$ can be viewed as an element $y \in \Hom_{\bbL_{\star \hspace{-0.1em}}[G]}(\overline{W},W)$ for which $xy = yx = 1$. For the third point note that $W|_{\bbK[G]} \cong U^{\oplus n}$ for $n \in \{1,2\}$ if $U \cong U \otimes \pi$ and $W|_{\bbK[G]} \cong U \oplus (U \otimes \pi)$ if $U \not\cong U \otimes \pi$ by Proposition \ref{prop:algindressummary}. Then as a $\bK$-algebra, $D(W) = \End_{\bbK[G]}(W)$ is $M_n(\End_{\bbK[G]}(U))$ in the first case and $\End_{\bbK[G]}(U) \times \End_{\bbK[G]}(U \otimes \pi)$ in the latter. The final point follows from Lemma \ref{lem:centreofC2graddiv}.
\end{proof}

Now we want to describe $\End_{\bR[H]}(W)$ in terms of $D(W)$. For this, we introduce the following.

\begin{defn}
For $A$ a $C_2$-graded algebra over $\bK$ and $\alpha \in \bL \setminus \bK$, we set
\[
A_{\alpha} \coloneqq A \, \hat{\otimes} \, \bK[\alpha],
\]
where $\bK[\alpha] = \bL$ as a $\bK$-algebra, with $C_2$-grading $\bK[\alpha]_0 = \bK$, $\bK[\alpha]_1 = \bK \cdot \alpha$.
\end{defn}

To establish some basic properties of $A_{\alpha}$, recall that for any $C_2$-graded algebra $A$ one can consider the \emph{main automorphism} of $A$, $t \colon A \rightarrow A$, which acts as the identity on $A_0$ and multiplication by $-1$ on $A_1$ \cite[Def.\ IV.3.7]{LAM}. When $A = D(W)$, this agrees with $t \colon D(W) \rightarrow D(W)$ as defined above.

\begin{lemma}\label{lem:Aalphaproperties}
    Suppose that $A$ is a $C_2$-graded algebra over $\bK$ and $\alpha \in \bL \setminus \bK$. Then:
    \begin{itemize}
        \item For a formal symbol $\underline{\alpha}$, the $\bK$-linear identification 
        \[
            \phi \colon A \oplus A \cdot \underline{\alpha} \xrightarrow{\sim} A_{\alpha}, \qquad \phi(a + b \cdot \underline{\alpha}) = a \otimes 1 + b \otimes \alpha
        \] 
        induces on $A \oplus A \cdot \underline{\alpha}$ multiplication determined by
\[
\underline{\alpha}^2 = \alpha^2, \qquad \underline{\alpha} x = t(x) \underline{\alpha}
\]
for $x \in A$, and so explicitly given by, for $x,y, u ,v \in A$,
        \[
        (x + y \cdot \underline{\alpha}) * (u + v \cdot \underline{\alpha}) = (xu + \alpha^2 y t(v)) + (xv + y t(u)) \cdot \underline{\alpha}.
        \]
        The induced $C_2$-grading is 
        \[
        (A_{\alpha})_0 = A_0 \oplus A_1 \cdot \underline{\alpha}, \qquad \text{and} \qquad (A_{\alpha})_1 = A_1 \oplus A_0 \cdot \underline{\alpha}.
        \]
        \item If $A$ is $C_2$-graded simple, then $A_{\alpha}$ is $C_2$-graded simple.
        \item If $\beta \in \bL \setminus \bK$ with $\langle \alpha \rangle_{\bbK} = \langle \beta \rangle_{\bbK}$, then $A_{\alpha} \cong B_{\alpha}$ as $C_2$-graded algebras.
        \item The inclusion $A \hookrightarrow A_{\alpha}$ induces identifications $\widehat{Z}(A) = \widehat{Z}(A_{\alpha})$ and
        \[
            Z(A_{\alpha}) = \widehat{Z}(A)_0 \oplus S \cdot \underline{\alpha}, \text{ where } S = \{x \in A_0 \mid x t(a) = ax \text{ for all } a \in A\}.
        \]
        When $A = D$ is a $C_2$-graded division algebra, then $\bF \coloneqq\widehat{Z}(D)$ is a field and this simplifies to
        \[
        Z(D_{\alpha}) = \twopartdef{\bF \oplus \bF u \cdot \underline{\alpha}}{\text{there exists $u \in D_0 \setminus 0$ with $u^{-1}xu = t(x)$ for any $x \in D$,}}{\bF}{\text{otherwise.}}
        \]
        The first case occurs if and only if $D$ is simple as an $\bF$-algebra and $Z(D) = \bF$, and the second case occurs if and only if $D_{\alpha}$ is simple as an $\bF$-algebra and has centre $\bF$.
    \end{itemize}
\end{lemma}

\begin{proof}
    The first point follows by direct computation. The second follows from the fact that the grading on $\bK[\alpha]$ only depends on $\langle \alpha \rangle_{\bbK}$. If $A$ is $C_2$-graded simple, then $A_{\alpha}$ is $C_2$-graded simple, and $A \hookrightarrow A_{\alpha}$ induces an identification $\widehat{Z}(A) = \widehat{Z}(A_{\alpha})$, both of which follow from \cite[Thm.\ IV.2.3]{LAM}, as $\bK[\alpha]$ is $C_2$-graded simple with graded centre $\bK$. It remains only to compute $Z(A_{\alpha})$, which, being graded, is determined by $Z(A_{\alpha})_0$ and $Z(A_{\alpha})_1$. We have that $Z(A_{\alpha})_0 = \widehat{Z}(A_{\alpha})_0 = \widehat{Z}(A)_0$.

    For $x \in A$, we have that $x \underline{\alpha} \in Z(A_{\alpha})_1$ if and only if for any $a \in A$, $x \underline{\alpha}$ commutes with $a$ and $a \underline{\alpha}$, or equivalently $x t(a) = ax$ and $x t(a) = a t(x)$ respectively. Taking $a = 1$, $t(x) = x$, and thus
    \[
    \{ x \in A \mid  x \underline{\alpha} \in Z(A_{\alpha})_1 \} = \{x \in A_0 \mid x t(a) = ax \text{ for all } a \in A\}.
    \]

    When $A = D$ is a $C_2$-graded division algebra, $\widehat{Z}(D) = \widehat{Z}(D)_0$ by Lemma \ref{lem:centreofC2graddiv}, and for $Z(A_{\alpha})_1$, if there is such a non-zero element $u$, then right-multiplication by $u$ defines an isomorphism
    \[
    Z(D) \cap D_0 = \widehat{Z}(D) \xrightarrow{\sim} \{x \in D_0 \mid x t(y) = yx\},
    \]
    from which we yield the required description. The final statement is \cite[Thm.\ IV.3.4]{LAM}.
\end{proof}

The following is the main result of this section, which allows us to relate $D(W)$ with $\End_{\bK[H]}(W)$. We consider $\End_{\bK[H]}(W)$ with $C_2$-grading induced by the order two automorphism
\[
t \coloneqq t_{\alpha,g} \colon \End_{\bbK[H]}(W) \rightarrow \End_{\bbK[H]}(W), \qquad t(f) = \rho(g)^{-1} \circ m_{\alpha^{-1}} \circ f \circ m_{\alpha} \circ \rho(g)
\]
for any $g \in G \setminus H$ and $\alpha \in \bL \setminus \bK$ with $\overline{\alpha} = - \alpha$, where $m_{\alpha} \colon W \rightarrow W$ denotes left-multiplication by $\alpha$. This is independent of the choice of $g$ and $\alpha$.

\begin{prop}\label{prop:recoverKHendo}
    Let $W \in \Irr_{\bbL}^{\hspace{0.1em} \star}(G)$ and $\alpha \in \bL \setminus \bK$ with $\overline{\alpha} = - \alpha$. Then the natural map
    \[
    \phi_{\alpha} \colon D(W)_{\alpha} \rightarrow \End_{\bbK[H]}(W), \qquad \phi_{\alpha}(f \otimes \lambda) = f \circ m_{\lambda},
    \]
    where $m_{\lambda}$ denotes left-multiplication by $\lambda \in \bK[\alpha]$, is an isomorphism of $C_2$-graded $\bK$-algebras.
\end{prop}

\begin{proof}
    With respect to the $C_2$-grading on $\End_{\bbK[H]}(W)$, both $D(W)$ and $\bK[\alpha]$ are $C_2$-graded subalgebras of $\End_{\bbK[H]}(W)$, where we view $\bK[\alpha] \hookrightarrow \End_{\bbK[H]}(W)$ via $\lambda \mapsto m_{\lambda}$. By definition, 
    \[
        \widehat{C}_{\End_{\bbK[H]}(W)}(\bK[\alpha]) = \{ f_0 + f_1 \in \End_{\bbK[H]}(W) \mid m_{\alpha^{-1}} \circ f_0 \circ m_{\alpha} = f_0, m_{\alpha^{-1}} \circ f_1 \circ m_{\alpha} = -f_1\},
    \]
    which, as $t(f_0) = f_0$ and $t(f_1) = -f_1$, are exactly the condition that both $f_0$ and $f_1$ commute with $\rho(g)$. As $f_0,f_1 \in \End_{\bbK[H]}(W)$ and thus commute with all $\rho(h)$ for $h \in H$, we thus obtain that 
    \[
    \widehat{C}_{\End_{\bbK[H]}(W)}(\bK[\alpha]) = \End_{\bbK[G]}(W) = D(W),
    \]
    from which the result follows by Proposition \ref{prop:graddoublecent}.
\end{proof}

The key point is that the isomorphism type of $D(W)$ as a $C_2$-graded division algebra determines the centre of the endomorphism algebra of $X$, as described in the following.

\begin{cor}\label{cor:Dalphaisom}
If $W \in \Irr_{\bbL}^{\hspace{0.1em} \star}(G)$, and $D \cong D(W)$ as $C_2$-graded $\bK$-algebras, then for any $\alpha \in \bL \setminus \bK$ with $\overline{\alpha} = - \alpha$,
\[
    D_{\alpha} \cong \End_{\bK[H]}(W)
\]
as $\bK$-algebras. In particular, $D_{\alpha}$ is simple as a $\bK$-algebra if and only if $X \cong {}^{\omega}X$, and
\[
Z(D_{\alpha}) \cong Z(\End_{\bbK[H]}(W)) \cong \twopartdef{Z(\End_{\bbK[H]}(X))}{X \cong {}^{\omega}X}{Z(\End_{\bbK[H]}(X)) \times Z(\End_{\bbK[H]}(X))}{X \not\cong {}^{\omega}X}
\]
as $\bK$-algebras.
\end{cor}

\begin{remark}
We note that, in practice, $Z(D_{\alpha})$ can be described using Lemma \ref{lem:Aalphaproperties}.
\end{remark}

\begin{proof}
    After Proposition \ref{prop:recoverKHendo}, all that remains to show is the final claim regarding $X$. We have that $W|_{k[H]} \cong X^{\oplus n}$ for some $n \geq 1$ if $X \cong {}^{\omega}X$ and $W|_{k[H]} \cong X^{\oplus n} \oplus ({}^{\omega}X)^{\oplus n}$ for some $n \geq 1$ if $X \not\cong {}^{\omega}X$ (which can be seen considering the restriction possibilities for simple modules from $\Rep_{\bbL}^{\hspace{0.1em} \star}(G)$ to $\Rep_{\bbK}(G)$, and $\Rep_{\bbK}(G)$ to $\Rep_{\bbK}(H)$). The claim then follows taking endomorphism rings.
\end{proof}

\begin{eg}
The extra restrictions on the type that Corollary \ref{cor:Dalphaisom} do not seem enough to rule out Case 21 and resolve Problem \ref{prob:case21}. However, for a given extension $\bL / \bK$, Corollary \ref{cor:Dalphaisom} does place significant restrictions on the type: for example, see Section \ref{sect:recoverDyson10} for the case $\bC / \bR$.
\end{eg}

Now we turn to $A(W)$, which we define in terms of blocks. Fix $\alpha \in \bL \setminus \bK$ with $\overline{\alpha} = - \alpha$. Conjugation by $\alpha$ defines an order-two automorphism of $\bK[G]$
\[
\phi \colon \bK[G] \rightarrow \bK[G], \qquad \phi(x) = \alpha x \alpha^{-1},
\]
which is independent of the choice of $\alpha$ and induces an order-two automorphism 
\[
\hat{\phi} \colon \widehat{\bK[G]} \rightarrow \widehat{\bK[G]}, \qquad \hat{\phi}((x_M)_M) = (x_{\phi_* N})_N.
\]
Concretely, $\phi(h) = h$ for each $h \in H$ and $\phi(g) = - g$ for each $g \in G \setminus H$, and thus $\phi_* N = N \otimes \pi$ is nothing but the twist of $N$ by the sign representation $\pi$ of $G$. Because $\phi$ and $\widehat{\phi}$ have order two and $\charfield(\bK) \neq 2$, the $\pm 1$ eigenspaces of $\phi$ and $\hat{\phi}$ define a $C_2$-gradings on $\bK[G]$ and $\widehat{\bK[G]}$ respectively.

It is clear that 
\[
\bK[G]_0 = \bK[H], \qquad \bK[G]_0 = \bK[G \setminus H].
\]
On the other hand, we need to use Lemma \ref{lem:gradedzeropartcompletion} to show the corresponding result for $\widehat{\bK[G]}$:

\begin{cor}\label{cor:degzeropartofcompletion}
    The inclusion
    \[
    \widehat{\iota}_{\bbK[G]} \colon \widehat{\bK[H]} \hookrightarrow \widehat{\bK[G]}
    \]
    induces an identification
    \[
    \widehat{\iota}_{\bbK[G]}(\widehat{\bK[H]}) = \widehat{\bK[G]}_0.
    \]
\end{cor}

For any central primitive idempotent $e \in \widehat{\bK[H]}$, we can consider the idempotent $f$ of $\widehat{\bK[H]}$ from Definition \ref{def:ideminbigger}. By definition, the idempotent $\iota_{\bbK[G]}(f)$ is central in $\widehat{\bK[G]}$, and therefore lies in $\widehat{\bK[G]}_0$, and thus the following is a well-defined $C_2$-grading.

\begin{defn}
    For $W \in \Irr_{\bbL}^{\hspace{0.1em} \star}(G)$, $A(W) \coloneqq \widehat{\iota}_{\bbK[G]}(f) \cdot \widehat{\bK[G]}$, with $C_2$-grading
    \[
        A(W)_0 = \widehat{\iota}_{\bbK[G]}(f) \cdot \widehat{\bK[G]}_0, \qquad A(W)_1 = \widehat{\iota}_{\bbK[G]}(f) \cdot \widehat{\bK[G]}_1.
    \]
\end{defn}

\begin{lemma}\label{lem:basicpropertiesAW}
    For $W \in \Irr_{\bbL}^{\hspace{0.1em} \star}(G)$, $A(W)$ is simple and non-trivially graded as a $C_2$-graded algebra over $\bK$, and
    \[
        A(W)_0 = \widehat{\iota}_{\bbK[G]}(f \cdot \widehat{\bK[H]}).
    \]
    Furthermore, $A(W)$ \emph{(}resp.\ $A(W)_0$\emph{)} is simple as a $\bK$-algebra if and only if $U \cong U \otimes \pi$ \emph{(}resp.\ $X \cong {}^{\omega}X$\emph{)}.
\end{lemma}

\begin{proof}
    The description of $A(W)_0$ follows directly from Corollary \ref{cor:degzeropartofcompletion}. To see that $A(W)$ is non-trivially graded, note that the image of any $g \in G \setminus H$ under the $C_2$-graded $\bK$-algebra homomorphism 
    \[
    \bK[G] \rightarrow \widehat{\bK[G]} \twoheadrightarrow \widehat{\iota}_{\bbK[G]}(f) \cdot \widehat{\bK[G]}
    \]
    is a unit, as $g$ is a unit of $\bK[G]$, and lies in $(\widehat{\iota}_{\bbK[G]}(f) \cdot \widehat{\bK[G]})_1$. 

    For the simplicity of $A(W)$, note that using Proposition \ref{prop:descofSblock} there are $\bK$-algebras isomorphisms
        \[
            A(W)_0 \cong \prod_{T \in C_2 \cdot X} M_{n_X}(\End_{\bbK[H]}(X)), \qquad A(W) \cong \prod_{T \in C_2 \cdot U} M_{n_U}(\End_{\bbK[G]}(U))
        \]
        for some $n_X, n_U \geq 0$. If $C_2 \cdot U = \{U\}$, then $A(W)$ is simple even as a $\bK$-algebra, so suppose instead that $U \not\cong U \otimes \pi$. From the four cases of Proposition \ref{prop:algindressummary} applied to the extension $\bK[H] \rightarrow \bK[G]$, the fact that $U \not\cong U \otimes \pi$ forces $X \cong {}^{\omega}X$, and thus $A(W)_0$ is simple as a $\bK$-algebra. Then $A(W)$ is $C_2$-graded simple by Lemma \ref{lem:invoddeltgradsimple}, using the invertible odd element of $A(W)$ described above.
\end{proof}

We can now describe, for $W \in \Irr_{\bbL}^{\hspace{0.1em} \star}(G)$, how $A(W)$ and $D(W)$ are related. To do this, we realise both $A(W)$ and $D(W)$ as $C_2$-graded subalgebras of $\End_{\bbK}(W)$.

Here, for any $W \in \Vect_{\bbL}$, we view $\End_{\bbK}(W)$ as a $C_2$-graded algebra with grading by linear and antilinear maps, as in the definition of $D(W)$. We write $\bK[\epsilon]$ for $\bK[\epsilon] / (\epsilon^2 - 1)$, with $C_2$-grading $\bK[\epsilon]_0 = \bK$, $\bK[\epsilon]_1 = \bK \cdot \epsilon$, and in the following view $M_{n}(\bK[\epsilon])$ with $C_2$-grading
\[
M_{n}(\bK[\epsilon])_0 = M_{n}(\bK), \qquad M_{n}(\bK[\epsilon])_1 = M_{n}(\bK \cdot \epsilon).
\]

We first describe the class of $\End_{\bbK}(W)$ in $\BrW(\bK)$.

\begin{lemma}
    If $W \in \Vect_{\bbL}$, then for any $\bL$-basis $e_1, ... ,e_n$ of $W$ the natural map
    \[
        \varphi \colon M_{n}(\bK[\epsilon]) \rightarrow \End_{\bbK}(W), \qquad \varphi((a_{ij}) + (b_{ij})\epsilon)\left(\sum_{k} \lambda_k e_k \right) \coloneqq \sum_{j,k} (a_{kj}\lambda_j + b_{kj} \overline{\lambda_j}) e_k
    \]
    is an isomorphism of $C_2$-graded $\bK$-algebras. In particular, $[\bK[\epsilon]] = [\End_{\bbK}(W)]$ in $\BrW(\bK)$.
\end{lemma}

\begin{proof}
    Because $\bL / \bK$ is Galois, the natural map $\bK[\Gal(\bL / \bK)] \rightarrow \End_{\bbK}(\bL)$ is an isomorphism. Then the map $\varphi$ is the composition of the canonical isomorphisms
    \[
    M_n(\bK[\epsilon]) \equiv M_n(\bK[\Gal(\bL / \bK)]) \xrightarrow{\sim} M_n(\End_{\bbK}(\bL)) \xrightarrow{\sim} \End_{\bbK}(\bL^n) \xrightarrow{\sim} \End_{\bbK}(W),
    \]
    and is easily seen to preserve the $C_2$-grading.
\end{proof}

The $C_2$-graded division $\bK$-algebra $D(W)$ is naturally a $C_2$-graded $\bK$-subalgebra of $\End_{\bbK}(W)$. To realise $A(W)$ as such a subalgebra requires a little more work. There is a natural $\bK$-algebra homomorphism
\[
 \widehat{\bK[G]} \rightarrow \End_{\bK}(W), \qquad (x_M)_M \mapsto x_{W|_{\bbK[G]}},
\]
which, as $U$ and $U \otimes \pi$ are the only constituents of $W|_{\bbK[G]}$, induces an injection of $\bK$-algebras
\[
\pi_W \colon A(W) = \widehat{\iota}_{\bbK[G]}(f) \cdot \widehat{\bK[G]} \hookrightarrow \End_{\bbK}(W).
\]

\begin{lemma}
    The inclusion
    \[
        \pi_W \colon A(W) \hookrightarrow \End_{\bbK}(W)
    \]
    preserves the $C_2$-grading, and thus is a homomorphism of $C_2$-graded $\bK$-algebras.
\end{lemma}

\begin{proof}
    There is similarly a natural $\bK$-algebra homomorphism
\[
 \widehat{\bK[H]} \rightarrow \End_{\bK}(W), \qquad (y_N)_N \mapsto y_{W|_{\bbK[H]}},
\]
which lands inside $\End_{\bbK}(W)_0 = \End_{\bbL}(W)$, because $m_i$ is a $\bK[H]$-linear endomorphism of $W|_{\bbK[H]}$, and so commutes with $y_{W|_{\bbK[H]}}$. This similarly induces an injection
\[
\eta_W \colon f \cdot \widehat{\bK[H]} \hookrightarrow \End_{\bbK}(W)_0
\]
because $X$ and ${}^{\omega} X$ are the only constituents of $W|_{\bbK[H]}$, which fits into the commutative diagram
% https://q.uiver.app/#q=WzAsMyxbMCwwLCJmIFxcY2RvdFxcd2lkZWhhdHtcXGJLW0hdfSJdLFsxLDEsIlxcRW5kX3tcXGJiS30oVykiXSxbMCwxLCJcXGhhdHtcXGlvdGF9X3tcXGJiS1tHXX0oZikgXFxjZG90IFxcd2lkZWhhdHtcXGJLW0ddfSJdLFswLDIsIlxcaGF0e1xcaW90YX1fe1xcYmJLW0ddfSIsMix7InN0eWxlIjp7InRhaWwiOnsibmFtZSI6Imhvb2siLCJzaWRlIjoidG9wIn19fV0sWzIsMSwiXFxwaV9XIiwwLHsic3R5bGUiOnsidGFpbCI6eyJuYW1lIjoiaG9vayIsInNpZGUiOiJ0b3AifX19XSxbMCwxLCJcXGV0YV9XIiwwLHsic3R5bGUiOnsidGFpbCI6eyJuYW1lIjoiaG9vayIsInNpZGUiOiJ0b3AifX19XV0=
\[\begin{tikzcd}
	{f \cdot\widehat{\bK[H]}} & \\
	{\hat{\iota}_{\bbK[G]}(f) \cdot \widehat{\bK[G]}} & {\End_{\bbK}(W)}
	\arrow["{\hat{\iota}_{\bbK[G]}}"', hook, from=1-1, to=2-1]
	\arrow["{\eta_W}", hook, from=1-1, to=2-2]
	\arrow["{\pi_W}", hook, from=2-1, to=2-2]
\end{tikzcd}\]
In particular, using Corollary \ref{cor:degzeropartofcompletion},
\[
\pi_W(A(W)_0) = \pi_W(\hat{\iota}_{\bbK[G]}(f \cdot\widehat{\bK[H]})) = \eta_W(f \cdot\widehat{\bK[H]}) \subset \End_{\bbK}(W)_0.
\]
For $A(W)_1$, the action map $\bK[G] \rightarrow \End_{\bbK}(W)$ is $C_2$-graded and is part of the commutative diagram
% https://q.uiver.app/#q=WzAsMyxbMCwwLCJcXGJLW0ddIl0sWzEsMSwiXFxFbmRfe1xcYmJLfShXKSJdLFswLDEsIlxcd2lkZWhhdHtcXGJLW0ddfSJdLFswLDIsIiIsMix7InN0eWxlIjp7InRhaWwiOnsibmFtZSI6Imhvb2siLCJzaWRlIjoidG9wIn19fV0sWzIsMV0sWzAsMV1d
\[\begin{tikzcd}
	{\bK[G]} & \\
	{\widehat{\bK[G]}} & {\End_{\bbK}(W)}
	\arrow[hook, from=1-1, to=2-1]
	\arrow[from=1-1, to=2-2]
	\arrow[from=2-1, to=2-2]
\end{tikzcd}\]
Thus any $g \in G \setminus H$ defines an element of $A(W)^{\times} \cap A(W)_1$ for which $\pi_W(g) \in \End_{\bbK}(W)_1$, so
\[
\pi_W(A(W)_1) = \pi_W(g \cdot A(W)_0) = \pi_W(g) \cdot \pi_W(A(W)_0) \subset \pi_W(g) \cdot \End_{\bbK}(W)_0 = \End_{\bbK}(W)_1
\]
as required.
\end{proof}

We want to relate $D(W)$ to the graded centraliser of the image of $A(W)$ in $\End_{\bbK}(W)$. 

\begin{defn}
For $W \in \Irr_{\bbL}^{\hspace{0.1em} \star}(G)$ we set
\[
\tilde{D}(W) \coloneqq \widehat{C}_{\End_{\bbK}(W)}(\pi_W(A(W))).
\]
\end{defn}

The $C_2$-graded $\bK$-subalgebras $D(W)$ and $\tilde{D}(W)$ are not equal:
\[
D(W)_0 = \tilde{D}(W)_0, \qquad D(W)_1 \cap \tilde{D}(W)_1 = 0.
\]
However, for any $\alpha \in \bL \setminus \bK$ with $\overline{\alpha} = - \alpha$, multiplication by $\alpha$ defines a $\bK$-linear isomorphism $D(W)_1 \xrightarrow{\sim} \tilde{D}(W)_1$. The below shows that if additionally $\alpha^2 = -1$, this extends to a $C_2$-graded $\bK$-algebra isomorphism $\psi_W \colon D(W) \rightarrow \tilde{D}(W)$, and thus allows us to relate $A(W)$ and $D(W)$.

\begin{prop}\label{prop:relateAWDW}
Suppose that $\bL = \bK(i)$ where $i^2 = -1$, and $W \in \Irr_{\bbL}^{\hspace{0.1em} \star}(G)$. Then the map
\[
   \psi_W \colon D(W) \rightarrow \tilde{D}(W), \qquad f_0 + f_1 \mapsto f_0 + if_1
\]
is an isomorphism of $C_2$-graded $\bK$-algebras.

Consequently, $\widehat{Z}(A(W)) = \bK$ if and only if $\widehat{Z}(D(W)) = \bK$, in which case the natural map
\[
\pi_W \otimes \psi_W \colon A(W) \, \hat{\otimes} \, D(W) \rightarrow \End_{\bbK}(W)
\]
is an isomorphism of $C_2$-graded $\bK$-algebras and thus an equality
\[
[A(W)] \cdot [D(W)] = [\bK[\epsilon]]
\]
in $\BrW(\bK)$.
\end{prop}

\begin{proof}
    The image $\pi_W(A(W))$ is the same as the image of the action map $\bK[G] \rightarrow \End_{\bbK}(W)$, both being equal to $\End_{\End_{\bbK[G]}(W)}(W)$. In particular, $\tilde{D}(W)_0 = \End_{\bbK[G]}(W)_0 = D(W)_0$ and 
    \[
        \tilde{D}(W)_1 = \{f \in \End_{\bbK}(W) \mid  f \circ m_i = - m_i \circ f, \rho(g) \circ f = \sigma(g) f \circ \rho(g) \text{ for all } g\in G \}
    \]
    where $\sigma \colon G \twoheadrightarrow \{\pm 1\}$ is the structure map of $G$. It is clear that
    \[
    D(W)_1 \rightarrow \tilde{D}(W)_1, \qquad f_1 \mapsto if_1
    \]
    is a $\bK$-linear isomorphism. The condition that $i^2 = -1$ is used to show that $\psi_W$ is multiplicative:
    \begin{align*}
        \psi_W(x_0 + x_1)\psi_W(y_0 + y_1) &= (x_0 + i x_1)(y_0 + i y_1), \\
        &= x_0y_0 + (ix_1)(i y_1) + x_0 (iy_1) + ix_1y_0, \\
        &= x_0y_0 - i^2 x_1y_1 + i x_0y_1 + ix_1y_0, \\
        &= (x_0y_0 + x_1y_1) + i(x_0y_1 + x_1 y_0), \\
        &= \psi_W((x_0 + x_1)(y_0 + y_1)),
    \end{align*}
    so $\psi_W$ is an isomorphism of $C_2$-graded $\bK$-algebras. Proposition \ref{prop:graddoublecent} shows the remaining claims.
\end{proof}

\section{Examples}

\subsection{The Split Case}

Suppose that $G = H \times C_2$, and $V \in \Irr_{\bL}(H)$. Then $U \not\cong U \otimes \pi$, and consulting Table \ref{tab:24types} and using Theorem \ref{thm:24ways}, we thus have the following.

\begin{cor}
    Suppose that $G = H \times C_2$, and $V \in \Irr_{\bL}(H)$. Then $\fT(V) \in \{(1), (10), (16)\}$.
\end{cor}

These correspond to the semilinear type of $V$ being A, B or C respectively, as described in Theorem \ref{thm:threeoptions}. Note also that the type of $V$ for the extension $\bK[H] \subset \bL[H]$ is the same as the type of $V$ for the extension $\bL[H] \subset \bL_{\star \hspace{-0.1em}}[G]$, which is consistent with the equivalence $\Rep_{\bbL}^{\hspace{0.1em} \star}(G) \cong \Rep_{\bK}(H)$ of \cite[\S 6]{TAYSchur}.

\subsection{Symmetric Groups}

Suppose that $G = S_n$ and $H = A_n$ for some $n \geq 3$ (to ensure that $H$ is index two in $G$). Suppose that $\charfield(\bK) = 0$ or $\charfield(\bK) > n$, and let $V \in \Irr_{\bbL}(A_n)$. The key property that restricts the type of $V$ is that every linear representation of $S_n$ over $\bK$ is absolutely irreducible. This forces $V$ to have type A for the extension $\bK[S_n] \leq \bL[S_n]$ and $\dim_{\bbL} \End_{\bbL[S_n]}(Z) = 1$. Consulting Tables \ref{tab:24types} and \ref{tab:24types6}, and using Theorem \ref{thm:24ways}, we thus immediately have the following.

\begin{cor}\label{cor:symgroups}
Suppose that $n \geq 3$, and $V \in \Irr_{\bbL}(A_n)$. Then $\fT(V) \in \{(1), (2), (3), (12)\}$.

In particular, the semilinear Schur index $m_{\bbK}^{\bbL}(V)$ of $V$ is $1$.
\end{cor}

In particular, as described in Section \ref{sect:semilinearschurindex}, this allows on to completely determine the semilinear representation theory of $S_n$ over $\bL$ from the character table of $A_n$ over $\bL$ together with its action of $C_2$ (cf.\ Section \ref{sect:C2actionIrrH}).

\section{Further Questions}\label{sect:24typesfurtherquestions}

We end this chapter with some further questions, in addition to Problem \ref{prob:case21}, on detailed aspects of the structure of the 24 types.

\subsection{Realisation for fixed extension}

In Section \ref{sect:realise24ways} we considered the realisability of each type for some quadratic extension and finite Real structure. The first question is to obtain a finer understanding once the field extension is fixed.

\begin{prob}
Suppose that $\bL / \bK$ is a given quadratic field extension. Which of the 24 types occur for some finite group $H$, Real structure $G$ on $H$, and irreducible representation $V \in \Irr_{\bbL}(H)$?
\end{prob}

For example, if $\bL / \bK$ is a number field, it is natural to ask which arithmetic properties of $\bL / \bK$ determine which types occur. The following shows that when $\bK$ is totally real, not all of the twelve types of Section \ref{sect:LsplitsHG} are realised.

\begin{eg}
Suppose that $\bL$ is a CM number field, $\bK$ is the index two totally real subfield, and $V \in \Irr_{\bbL}(H)$ is absolutely irreducible. Then the D-type of $V$ cannot be II.V.

Indeed, if the linear type of $V$ is (B, A, 1, 3(b)), then by Table \ref{tab:24types6}, $Z$ is also absolutely irreducible, and the pair $(Z_{\bC}, V_{\bC})$ correspond over $\bC$. We have that $\Stab_{K_4}(Z_{\bC}) = \Stab_{K_4}(Z) = \langle c \rangle$ and $\Stab_{K_4}(V_{\bC}) = \Stab_{K_4}(V) = K_4$, and therefore the linear type of $V_{\bC}$ is also (B, A, 1, 3(b)) from Table \ref{tab:24types4}. This forces $V_{\bC}$ to have type II, as we will see in Theorem \ref{thm:dyson10}, and therefore $V_{\bC}$ does not extend to a semilinear representation of $G$ over $\bC$. This means that $V$ also doesn't extend to a semilinear representation of $G$: any such extension $W$ would define an semilinear extension $W_{\bC}$ of $V_{\bC}$. This forces $m_{\bK}^{\bL}(V) = 1$, and therefore $V$ to have D-type II, not D-type II.V.
\end{eg}

\begin{eg}
Suppose that $H$ has exponent $e$ and $G$ has exponent $f$. Then the possible types for the cyclotomic extensions $\bQ(\zeta_e)_+ \leq \bQ(\zeta_e)$ and $\bQ(\zeta_f)_+ \leq \bQ(\zeta_f)$ are described by Theorem \ref{thm:15waysforcyclotomic}, which also provides an computable method to determine the type.
\end{eg}

\subsection{Realisation of Brauer-Wall classes}

It is well-known that not all classes of the Brauer group of a field $\bK$ are realised as endomorphism rings of irreducible representations of finite groups over $\bK$, such classes which are realised in the way defining the Schur subgroup of $\Br(\bK)$, which is typically not equal to all of $\Br(\bK)$.

On the other hand, in \cite[\S 16]{TAYSchur} it is shown that any class $[A] \in \Br(\bK)$ is realised as $[\End_{\bbL_{\sigma}[G]}(V)]$ for some finite Galois extension $\bL / \bK$, group $G$, surjective homomorphism $\sigma \colon G \rightarrow \Gal(\bL / \bK)$ and $V \in \Irr_{\bL}^{\sigma}(G)$. The group $G$ can furthermore be taken to be finite if $[A]$ has order coprime to the characteristic of $\bK$. 

When $\bL / \bK$ has degree two, we saw in Section \ref{sect:BrauerWall} a way to associate to any S-block an element
\[
[D(W)] \in \BrW(\bK)
\]
whenever $\widehat{Z}(D(W)) = \bK$ (for example, this happens whenever $Z(\End_{\bbL_{\star \hspace{-0.1em}}[G]}(W)) = \bK$). When $\bL / \bK$ is $\bC / \bR$, $\BrW(\bR) \cong \bZ / 8 \bZ$ (see Proposition \ref{prop:gradeddivalgoverR} below), and this association is exhaustive.

It is natural to ask if the same holds for any quadratic $\bL / \bK$, and if not to determine the image.

\begin{prob}
Suppose that $\bL / \bK$ is a given quadratic field extension. Which classes of $\BrW(\bK)$ are realised for S-blocks of finite Real structures $H \leq G$ over $\bL / \bK$?
\end{prob}

\subsection{Prime degree classification}

Let $p$ be a prime number. Much of the theory of Section \ref{sect:fourfoldway} and this entire chapter goes through for prime degree extensions $\bL / \bK$, provided that $\bK$ contains a primitive $p$th root of $1$. 

\begin{prob}
Generalise the classification into 24 types to arbitrary prime degree $\bL / \bK$.
\end{prob}

More generally, one could try to obtain a similar classification into types of semilinear representations for $\Gal(\bL/ \bK)$ a fixed group, the case here being for the group $C_2$.

